\chapter{Radiation, Magnitudes \& Spectra}

\section*{Theory problems}

\probhead{6.1}{}{2007}{Chiang Mai, Thailand}{TH 1.10}{2}
% source id: 2007_TH_1.10   tag: distance from flux and luminosity
A K star on the Main Sequence has a luminosity of $0.4\,L_\odot$. This star
is observed to have a flux of $6.23 \times 10^{-14}\,\mathrm{W\,m^{-2}}$.
What is the distance to this star? You may neglect the atmospheric effect.

\probhead{6.2}{}{2007}{Chiang Mai, Thailand}{TH 1.11}{2}
% source id: 2007_TH_1.11   tag: supernova distance from flux comparison
A supernova shines with a luminosity $10^{10}$ times that of the Sun. If
such a supernova appears in our sky as bright as the Sun, how far away from
us must it be located?

\probhead{6.3}{}{2007}{Chiang Mai, Thailand}{TH 1.12}{2}
% source id: 2007_TH_1.12   tag: 21-cm Doppler velocity of gas cloud
The (spin-flip) transition of atomic hydrogen at rest generates the
electromagnetic wave of the frequency $\nu_0 = 1420.406$\,MHz. Such an
emission from a gas cloud near the galactic center is observed to have a
frequency $\nu = 1421.65$\,MHz. Calculate the velocity of the gas cloud. Is
it moving towards or away from the Earth?

\probhead{6.4}{}{2007}{Chiang Mai, Thailand}{TH 1.15}{2}
% source id: 2007_TH_1.15   tag: flux ratio from magnitude difference
The magnitude of the faintest star you can see with naked eyes is $m = 6$,
whereas that of the brightest star in the sky is $m = -1.5$. What is the
energy-flux ratio of the faintest to that of the brightest?

\probhead{6.5}{}{2007}{Chiang Mai, Thailand}{TH 1.8}{2}
% source id: 2007_TH_1.8   tag: Wien law peak wavelength 4000 K
At what wavelength does a star with the surface temperature of 4000\,K emit
most intensely?

\probhead{6.6}{}{2008}{Bandung, Indonesia}{TH L2}{100}
% source id: 2008_TH_L2   tag: UBV extinction, intrinsic colors, Teff
A $UBV$ photometric ($UBV$ Johnson's) observation of a star gives
$U = 8.15$, $B = 8.50$, and $V = 8.14$. Based on the spectral class, one gets
the intrinsic color $(U-B)_{\mathrm{o}} = -0.45$. If the star is known to
have radius of $2.3\,R_\odot$, absolute bolometric magnitude of $-0.25$, and
bolometric correction ($BC$) of $-0.15$, determine:
\begin{description}
\item[a.] the intrinsic magnitudes $U$, $B$, and $V$ of the star (take, for
  the typical interstellar matters, the ratio of total to selective
  extinction (color excess) $R_V = 3.2$),
\item[b.] the effective temperature of the star,
\item[c.] the distance to the star in pc.
\end{description}
Note: The relation between color excess of $U-B$ and of $B-V$ is
$E(U-B) = 0.72\,E(B-V)$.

\probhead{6.7}{}{2009}{Tehran, Iran}{TH P2}{10}
% source id: 2009_TH_P2   tag: photon rate into eye from m=6 star
Estimate the number of photons per second that arrive on our eye at
$\lambda = 550$\,nm (V-band) from a G2 main sequence star with apparent
magnitude of $m = 6$ (the threshold of naked eye visibility). Assume the eye
pupil diameter is 6\,mm and all the radition % sic: "radition" for "radiation" in the source
from this star is in
$\lambda = 550$\,nm.

\probhead{6.8}{}{2010}{Beijing, China}{TH S1}{10}
% source id: 2010_TH_S1   tag: combined magnitude of two stars
In a binary system, the apparent magnitude of the primary star is 1.0 and
that of the secondary star is 2.0. Find the maximum combined magnitude of
this system.

\probhead{6.9}{}{2011}{Katowice, Poland}{TH S8}{10}
% source id: 2011_TH_S8   tag: dust grain radiation pressure vs gravity
Assuming that dust grains are black bodies, determine the diameter of a
spherical dust grain which can remain at 1\,AU from the Sun in equilibrium
between the radiation pressure and gravitational attraction of the Sun. Take
the density of the dust grain to be
$\varrho = 10^{3}\,\mathrm{kg\,m^{-3}}$.

\probhead{6.10}{}{2011}{Katowice, Poland}{TH S9}{10}
% source id: 2011_TH_S9   tag: sky covering fraction, surface brightness
Interstellar distances are large compared to the sizes of stars. Thus,
stellar clusters and galaxies which do not contain diffuse matter essentially
do not obscure objects behind them. Estimate what proportion of the sky is
obscured by stars when we look in the direction of a galaxy of surface
brightness $\mu = 18.0\,\mathrm{mag\,arcsec^{-2}}$. Assume that the galaxy
consists of stars similar to the Sun.

\probhead{6.11}{}{2012}{Rio de Janeiro, Brazil}{TH L1}{}
% source id: 2012_TH_L1   tag: globular cluster combined magnitudes
An astronomer on Earth observes a globular cluster, which has an angular
diameter $\alpha$ and contains $N$ stars, each one with the same absolute
magnitude $M_0$, and is at a distance $D$ from the Earth. A biologist is in
the center of that cluster.
\begin{description}
\item[1.1.] What is the difference between the combined visual magnitudes of
  all stars observed by the astronomer and the biologist. \textbf{Consider
  that the spatial distribution of stars in the cluster is perfectly
  homogeneous and the biologist is measuring the combined magnitude of the
  entire cluster.}
\item[1.2.] What is the diameter of the astronomer's telescope, considering
  he wants to visualize the cluster with the same brightness that the
  biologist sees?
\item[1.3.] What would be the difference between the visual magnitudes
  observed by the two scientists, if the diameter of the field of view of
  the biologist is also $\alpha$.
\end{description}

\probhead{6.12}{}{2012}{Rio de Janeiro, Brazil}{TH S9}{}
% source id: 2012_TH_S9   tag: interstellar+nebular extinction of WDs
An old planetary nebula, with a white dwarf (WD) in its center, is located
50\,pc away from Earth. Exactly in the same direction, but behind the
nebula, lies another WD, identical to the first, but located at 150\,pc from
the Earth. Consider that the two WDs have absolute bolometric magnitude
$+14.2$ and intrinsic color indexes $B-V = 0.300$ and $U-V = 0.330$.
Extinction occurs in the interstellar medium and in the planetary nebula.

When we measure the color indices for the closer WD (the one who lies at the
center of the nebula), we find the values $B-V = 0.327$ and $U-B = 0.038$.
In this part of the Galaxy, the interstellar extinction rates are 1.50, 1.23
and 1.00 magnitudes per kiloparsec for the filters $U$, $B$ and $V$,
respectively. Calculate the color indices as they would be measured for the
second star.

\probhead{6.13}{}{2013}{Volos, Greece}{TH S13}{10}
% source id: 2013_TH_S13   tag: fog optical depth from Sun dimming
Recently in London, because of a very thick layer of fog, the visual
magnitude of the Sun, became equal to the (usual -- as observed during
cloudless nights) magnitude of the full Moon. Assuming that the reduction of
the intensity of light due to the fog is given by an exponential equation,
calculate the exponential coefficient, $\tau$, which is usually called
\emph{optical depth}.

\probhead{6.14}{}{2013}{Volos, Greece}{TH S8}{10}
% source id: 2013_TH_S8   tag: decompose combined magnitude of pair
A binary system of stars consists of star ($a$) and star ($b$) with
brightness ratio 2. The binary system is difficult to resolve and is
observed from the Earth as \underline{one star} of $5^{\mathrm{th}}$
magnitude. Find the apparent magnitude of each of the two stars
($m_a$, $m_b$).

\probhead{6.15}{Space-ship orbiting the Sun}{2014}{Suceava, Romania}{TH P11}{}
% source id: 2014_TH_P11   tag: Sun magnitude from ship equilibrium temperature
A spherical space-ship orbits the Sun on a circular orbit, and spin around
one of its axes. The temperature on the exterior surface of the ship is
$T_N$. Find out the apparent magnitude of the Sun and the angular diameter
of the Sun as seen by the astronaut on board of the space-ship. The
following values are known: $T_S$ -- the effective temperature of the Sun;
$R_S$ -- the radius of the Sun; $d_0$ -- the Earth--Sun distance; $m_0$ --
apparent magnitude of Sun measured from Earth; $R_N$ -- the radius of the
space-ship.

\probhead{6.16}{The effective temperature on the surface of a star}{2014}{Suceava, Romania}{TH P7}{}
% source id: 2014_TH_P7   tag: Teff from two spectral intensities
A star emits radiation with wavelength values in a narrow range
$\Delta\lambda \ll \lambda$, i.e. the wavelength have values between
$\lambda$ and $\lambda + \Delta\lambda$. % sic: the source prints "between λ and λ and λ + Δλ" (duplicated "λ and"); normalized here
According to Planck's relationship
(for an absolute black body), the following relation define, the energy
emitted by star in the unit of time, through the unit of area of its
surface, per length-unit of the wavelength range:
% sic: the exponent below is printed as hc/k\lambda c in the original paper
% (physically it should be hc/\lambda kT); transcribed as printed.
\[ r = \frac{2\pi h c^2}{\lambda^5\left(e^{hc/k\lambda c} - 1\right)}. \]

The spectral intensities of two radiations with wavelengths $\lambda_1$ and
respectively $\lambda_2$, both in the range $\Delta\lambda$, measured on
Earth are $I_1(\lambda_1)$ and, respectively $I_2(\lambda_2)$.
\begin{description}
\item[a.] Establish the equation which, in a general case, allows
  determining the effective temperature on the surface of the star using
  only spectral measurements.
\item[b.] Find out the approximate value of the effective temperature on the
  star surface if $hc \gg \lambda k T$.
\item[c.] Find out the relation between wavelength $\lambda_1$ and
  $\lambda_2$, if $I_1(\lambda_1) = 2 I_2(\lambda_2)$, when
  $hc \ll \lambda k T$.
\end{description}
You know: $h$ -- Planck's constant; $k$ -- Boltzmann's constant; $c$ --
speed of light in vacuum.

\probhead{6.17}{Gradient temperatures}{2014}{Suceava, Romania}{TH P8}{}
% source id: 2014_TH_P8   tag: temperature from two-line magnitude differences
The spectra of two stars with different temperatures $T_1$ and respectively
$T_2$ were compared. In the spectrum of each star, two very close spectral
lines corresponding to the wavelength with values $\lambda_1$ and respectively
$\lambda_2$ were found. For each line of this spectral lines, the difference
between the corresponding visual apparent magnitudes of the stars are
$\Delta m_{\lambda_1} = m_{1,\lambda_1} - m_{2,\lambda_1}$ and
$\Delta m_{\lambda_2} = m_{1,\lambda_2} - m_{2,\lambda_2}$.
$m_{1,\lambda_1}$ is the apparent magnitude of the star 1 for the wavelength
$\lambda_1$, $m_{1,\lambda_2}$ is the apparent magnitude of the star 1 for the
wavelength $\lambda_2$, $m_{2,\lambda_1}$ is the apparent magnitude of the
star 2 for the wavelength $\lambda_1$, $m_{2,\lambda_2}$ is the apparent
magnitude of the star 2 for the wavelength $\lambda_2$.

\emph{Determine} the temperature $T_1$ of one of the two stars, if the
temperature $T_2$ of the other star is already known, by using the Planck
expression of black body radiation:
\[ r(\lambda) = \frac{2\pi h c^2}{\lambda^5}
   \left( e^{hc/k\lambda c} - 1 \right)^{-1}, \]
% sic: the original paper prints the exponent as hc/k\lambda c (presumably a
% typo for hc/k\lambda T; cf. the closing condition below).
where: $h$ --- Planck's constant; $k$ --- Boltzmann's constant; $c$ --- speed
of light in vacuum. You will consider that $hc \gg k\lambda T$.

\probhead{6.18}{Pressure of light}{2014}{Suceava, Romania}{TH P9}{}
% source id: 2014_TH_P9   tag: dust grain radiation pressure vs gravity
One particle of star dust is in static equilibrium at a certain distance from
Sun. Assuming that the particle is spherical and its density is $\rho$,
calculate the diameter of the particle.

The following assumption may be useful for solving the problem: the pressure
of electromagnetic radiation is equal with the volume density of the
electromagnetic radiations.

\probhead{6.19}{}{2015}{Magelang, Indonesia}{TH S6}{}
% source id: 2015_TH_S6   tag: zenith optical depth from attenuation
At the start of every observation, a radio telescope is pointed at a
point-source calibrator that has a known flux density of 21.86\,Jy outside the
Earth's atmosphere. However, on a certain date, the measured flux density of
the calibrator source was 14.27\,Jy. If the calibrator source was at an
altitude of 35 degrees, estimate the zenith atmospheric optical depth,
$\tau_z$.

\probhead{6.20}{}{2015}{Magelang, Indonesia}{TH S8}{}
% source id: 2015_TH_S8   tag: brightness temperature of radio burst
A strong continuum radio signal from a celestial body has been observed as a
burst with a very short duration of 700\,$\mu$s. The observed flux density at
a frequency of 1660\,MHz is measured to be 0.35\,kJy. If the distance from the
source is known to be 2.3\,kpc, estimate the brightness temperature of this
source.

\probhead{6.21}{21-cm HI galaxy survey}{2017}{Phuket, Thailand}{TH T4}{10}
% source id: 2017_TH_T4   tag: HI flux vs detection limit, max redshift
A radio telescope is equipped with a receiver which can observe in a frequency
range from 1.32 to 1.52\,GHz. Its detection limit is 0.5\,mJy per beam for a
1-minute integration time. In a galaxy survey, the luminosity of the HI
spectral line of a typical target galaxy is $10^{28}$\,W with a linewidth of
1\,MHz. For a large beam, the HI emitting region from a far-away galaxy can be
approximated as a point source. The HI spin-flip spectral line has a
rest-frame frequency of 1.42\,GHz.

What is the highest redshift, $z$, of a typical HI galaxy that can be detected
by a survey carried out with this radio telescope, using 1-minute integration
time? You may assume in your calculation that the redshift is small and the
non-relativistic approximation can be used. Note that
$1\,\mathrm{Jy} = 10^{-26}\,\mathrm{W\,m^{-2}\,Hz^{-1}}$.

\probhead{6.22}{Supernova 1987A}{2017}{Phuket, Thailand}{TH T6}{15}
% source id: 2017_TH_T6   tag: magnitude decay, limiting magnitude
Supernova SN 1987A was at its brightest with an apparent magnitude of $+3$ on
about 15th May 1987 and then faded, finally becoming invisible to the naked
eye by 4th February 1988. It is assumed that brightness $B$ varied with time
$t$ as an exponential decline, $B = B_0 e^{-t/\tau}$, where $B_0$ and $\tau$
are constant. The maximum apparent magnitude which can be seen by the naked
eye is $+6$.

\begin{description}
\item[a)] Determine the value of $\tau$ in days. \hfill[5]
\item[b)] Find the last day that observers could have seen the supernova if
  they had a 6-inch (15.24-cm) telescope with transmission efficiency
  $T = 70\%$. Assume that the average diameter of the human pupil is
  0.6\,cm. \hfill[10]
\end{description}

\probhead{6.23}{Observe the Sun with FAST}{2018}{Beijing, China}{TH T6}{25}
% source id: 2018_TH_T6   tag: Rayleigh-Jeans solar radio flux
The Five-hundred-meter Aperture Spherical radio Telescope (FAST) is a
single-dish radio telescope located in Guizhou Province, China. The physical
diameter of the dish is 500\,m, but during observations, the effective
diameter of the collecting area is 300\,m.

Consider observations of the thermal radio emission from the photosphere of
the Sun at 3.0\,GHz with this telescope and a receiver with bandwidth
0.3\,GHz.

\begin{description}
\item[(a)] Calculate the total energy ($E_\odot$) that the receiver will
  collect during 1 hour of observation.
\item[(b)] Estimate the energy needed to turn over one page of your answer
  sheet ($E'$). Hint: the typical surface density of paper is
  80\,g\,m$^{-2}$.
\item[(c)] Which one is larger?
\end{description}

\probhead{6.24}{Cosmic Microwave Background Oven}{2019}{Keszthely, Hungary}{TH 5}{10}
% source id: 2019_TH_5   tag: blackbody CMB absorption, heating
Since the human body is made mostly of water, it is very efficient at
absorbing microwave photons. Assume that an astronaut's body is a perfect
spherical absorber with mass of $m = 60$\,kg, and its average density and heat
capacity are the same as for pure water, i.e.\
$\rho = 1000\,\mathrm{kg\,m^{-3}}$ and
$C = 4200\,\mathrm{J\,kg^{-1}\,K^{-1}}$.

\begin{description}
\item[a)] What is the approximate rate, in watts, at which an astronaut in
  intergalactic space would absorb radiative energy from the Cosmic Microwave
  Background (CMB)? The spectral energy distribution of CMB can be
  approximated by blackbody radiation of temperature
  $T_{\mathrm{CMB}} = 2.728$\,K. \hfill[5\,p]
\item[b)] Approximately how many CMB photons per second would the astronaut
  absorb? \hfill[3\,p]
\item[c)] Ignoring other energy inputs and outputs, how long would it take for
  the CMB to raise the astronaut's temperature by $\Delta T = 1$\,K?
  \hfill[2\,p]
\end{description}

\probhead{6.25}{Dyson Sphere}{2022}{Kutaisi, Georgia}{TH 11}{50}
% source id: 2022_TH_11   tag: blackbody equilibrium, radiation pressure
The Kardashev scale distinguishes three stages of evolution of civilizations
according to the criterion of access to and use of energy.

A type II civilization is capable of harnessing all the energy radiated by its
own star. Currently, we are a type zero civilization (we are not even
harnessing 100\% of the energy that reaches the Earth). One of the ways to
become a type II civilization is by building a Dyson Sphere. You can imagine
it as a sphere, built around the Sun, having the inner surface covered with
solar panels.

We assume that modern solar panels are used to build the sphere. First, let's
find out at what distance from the Sun should it be built. Emissivity of the
back side of solar panels is $\epsilon = 0.8$.

\begin{description}
\item[(a)] Solar panels absorb and transfer about $k = 30\%$ of the incident
  radiation into internal heat. Find the temperature of a Dyson Sphere of
  radius $R$. Express your answer in terms of $k$, $R$, $\epsilon$ and
  $L_\odot$ (solar luminosity). You may consider the Sun to be a black body.
  Ignore reflections from the solar panels. Ignore any possible effect of the
  energy not transferred to the internal heat of the solar panels or into
  electrical energy generated by the panels. \hfill(3\,pt)
\end{description}

Assume that the highest operational temperature for modern solar panels is
about $T_{max} \approx 104.5\,^\circ$C. After that, efficiency drops
significantly. To minimize the amount of material used, we should consider
building the sphere as small as possible.

\begin{description}
\item[(b)] Calculate the radius of the sphere for the panels to work properly.
  Does the Earth stay inside or outside the sphere? Write IN or OUT in the
  answer sheet. \hfill(4\,pt)
\item[(c)] Find the power harnessed by this Dyson Sphere, if the final power
  output of modern solar panels is about $\eta = 20\%$ of the incident energy.
  \hfill(2\,pt)
\item[(d)] Currently, the average power usage of the whole world is about 17
  Terawatts. If this Dyson sphere collects energy for one second, for how long
  could that meet our energy needs? \hfill(2\,pt)
\item[(e)] In the case where the Dyson sphere completely blocks out the rays
  of the Sun, the temperature on Earth will drop significantly. Calculate the
  change in the average temperature of the Earth in this condition, if current
  average temperature is about $15\,^\circ$C. Assume that Earth is also a
  black body. \hfill(5\,pt)
\item[(f)] Building a rigid spherical object of that size is nearly
  impossible. Another way of ``building'' the sphere is by sending individual
  panels to orbit around the sun (in different inclined orbits) at the radius
  $R$ found in part b. Calculate the period $T$ of any object orbiting the sun
  at that radius. \hfill(3\,pt)
\item[(g)] Assume that each solar panel is a thin sheet of silicon, having
  unit surface mass $\rho = 1\,\mathrm{kg/m^2}$. The radiation pressure from
  the Sun, might interfere with the orbit of the panel. Calculate the ratio
  $\alpha$ of the gravitational and photon forces for unit surface area of
  panels at distance $R$. You assume that all incident light is absorbed. Will
  this radiation pressure have any measurable effect? Write YES or NO in the
  Answer Sheet. \hfill(13\,pt)
\end{description}

Now assume that the Dyson Sphere is a rigid body rotating with the period
found in part `f' and having the radius found in part `b'.

\begin{description}
\item[(h)] A major threat to the Dyson Sphere would be an asteroid whose orbit
  crosses the surface of the sphere. One way to solve this problem is by
  removing panels from the path of the asteroid. Obviously the size of the
  hole through which asteroid should pass is much smaller than radius of the
  sphere.

  Astronomers discovered an asteroid on an orbit in the ecliptic plane in the
  same direction as the direction of rotation of the Dyson sphere. They
  calculated that this asteroid will enter the sphere on 14th of August and
  leave it on 20th of September. Calculate the angular distance between the
  two holes in the sphere needed to provide a safe passage for the asteroid.
  \hfill(13\,pt)

  The trajectory of this asteroid in heliocentric cylindrical coordinate
  system can be described as
  \[ r = \frac{a}{1 - \cos\theta} \]
  where $a = 1.00$\,au.
\item[(i)] In what range of wavelengths should we be searching for a Dyson
  sphere created by a type II civilization in a distant galaxy, if its
  distance from Earth is $d$ and the sphere can operate between temperatures
  $T_1$ and $T_2$ ($T_1 < T_2$). Assume only non-relativistic effects.
  \hfill(5\,pt)
\end{description}

\probhead{6.26}{Magnetic Field}{2023}{Katowice, Poland}{TH Q2}{5}
% source id: 2023_TH_Q2   tag: cyclotron line gives magnetic field
An emission line of wavelength $\lambda = 600$\,nm was observed in the
spectrum of a white dwarf. Assuming that it originates from the interaction of
a free non-relativistic electron with a magnetic field,

\begin{description}
\item[(a)] calculate the magnetic flux density of the field;
\item[(b)] estimate the wavelength of another spectral line, the discovery of
  which could confirm that the lines originate from particles of a plasma
  interacting with the magnetic field.
\end{description}

\probhead{6.27}{Balmer Decrement}{2025}{Mumbai, India}{TH T04}{15}
% source id: 2025_TH_T04   tag: Halpha/Hbeta ratio gives extinction
Consider a main sequence star surrounded by a nebula. The observed V-band
magnitude of the star is 11.315\,mag. The ionised region of nebula close to
the star emits H$\alpha$ and H$\beta$ lines; their wavelengths are
0.6563\,$\mu$m and 0.4861\,$\mu$m, respectively. The theoretically predicted
ratio of fluxes in H$\alpha$ to H$\beta$ lines is
$f_{\mathrm{H}\alpha}/f_{\mathrm{H}\beta} = 2.86$. However, when this
radiation passes through the outer portion of the cold dusty nebula, the
observed emission fluxes of H$\alpha$ and H$\beta$ lines are
$6.80 \times 10^{-15}\,\mathrm{W\,m^{-2}}$ and
$1.06 \times 10^{-15}\,\mathrm{W\,m^{-2}}$, respectively.

The extinction $A_\lambda$ is a function of wavelength and is expressed as
\[ A_\lambda = \kappa(\lambda)\,E(B-V). \]
Here, $\kappa(\lambda)$ is the extinction curve and $E(B-V)$ denotes the
colour excess in the filter bands B and V. The extinction curve (with
$\lambda$ in $\mu$m) is given as follows.
\[ \kappa(\lambda) =
   \begin{cases}
     2.659 \times \left( -1.857 + \dfrac{1.040}{\lambda} \right) + R_V, &
       0.63 \le \lambda \le 2.20 \\[1ex]
     2.659 \times \left( -2.156 + \dfrac{1.509}{\lambda}
       - \dfrac{0.198}{\lambda^2} + \dfrac{0.011}{\lambda^3} \right) + R_V, &
       0.12 \le \lambda < 0.63
   \end{cases} \]
where, $R_V = A_V/E(B-V) = 3.1$ is the ratio of total-to-selective extinction.

\begin{description}
\item[(T04.1)] Find the values of $\kappa(\mathrm{H}\alpha)$ and
  $\kappa(\mathrm{H}\beta)$. \hfill[3]
\item[(T04.2)] Find the value of the ratio of colour excess
  $\dfrac{E(\mathrm{H}\beta - \mathrm{H}\alpha)}{E(B-V)}$. \hfill[4]
\item[(T04.3)] Estimate the extinction due to nebula, $A_{\mathrm{H}\alpha}$
  and $A_{\mathrm{H}\beta}$, at H$\alpha$ and H$\beta$ wavelengths,
  respectively. \hfill[6]
\item[(T04.4)] Estimate the extinction of the nebula ($A_V$) and the apparent
  magnitude of the star in the V band, $m_{V0}$, in the absence of the nebula.
  \hfill[2]
\end{description}

\section*{Data-analysis problems}

\probhead{6.28}{}{2013}{Volos, Greece}{DA D3}{}
% source id: 2013_DA_D3   tag: visual magnitude estimation, Hyades chart
Figure 2 shows a photograph of the sky in the vicinity of the Hyades open
cluster. The V-filter in the Johnson's photometric system was used. Figure 3
is a chart of the region with known V-magnitudes ($m_V$) of several stars
(note that in order to avoid confusion with the stars, no decimal point is
used, i.e.\ a magnitude $m_V = 8.1$ is noted as ``81''). Hint: some of the
stars may not be in the chart.

\begin{description}
\item[(a)] Identify as many of the \emph{stars shown with a number and an
  arrow} in Figure 3 and mark them on Figure 2.
\item[(b)] Comparing the V-magnitudes of the known stars in Figure 2, estimate
  the V-magnitudes of the \emph{stars shown with a number and an arrow} in
  Figure 3.
\end{description}

\ioaafig{2013_DA_D3-f1.pdf}
% V filter) and Figure 3 (chart of the region with V-magnitudes of several
% stars) from the 2013 data-analysis paper are not in images/.

\ioaafig{2013_DA_D3-f2.pdf}

\probhead{6.29}{Observer on an extrasolar planet}{2014}{Suceava, Romania}{DA D3}{}
% source id: 2014_DA_D3   tag: magnitudes/distances recomputed elsewhere
The Sirius star, located in the constellation of Canis Major, is the brightest
star in the night sky of the Earth. What the observer's eye sees as a single
star is actually a binary star system. The high brightness of Sirius is a
consequence of two facts: its intrinsic luminosity and its proximity to the
Earth.

The Mizar multiple star system, in the constellation of Ursa Major, consists
of 4 stars seen along the same line of sight from the Earth. Some of these
stars form a gravitationally bound system.

Let's assume that an observer (observer A) is located on one of the planets of
the Sirius system. \emph{Determine}:

\begin{description}
\item[a)] The magnitude of the Sun as seen by observer A
  ($m_{\mathrm{Sun,Planet}}$).
\item[b)] The magnitude of Sirius star system as seen by the observer A.
  ($m_{\mathrm{SY,Planet}}$)
\item[c)] The combined intrinsic luminosity of the Mizar system,
  $L_{\mathrm{Mizar}}$;
\item[d)] the average distance between gravitationally bound stars of the
  Mizar system and Earth, $\ell$
  % sic: the symbol in the original prints as a stray grave accent; it is
  % most plausibly a script ell.
\item[e)] The geocentric angular distance between Mizar system and Sirius,
  $\Delta\theta$;
\item[f)] The physical distance between the gravitationally bound stars of
  the Mizar system and the observer A. ($d_{\mathrm{Mizar,Planet}}$)
\item[g)] The magnitude of the entire Mizar system as seen by the observer A.
  ($m_{\mathrm{Mizar,Planet}}$)
\end{description}

Also estimate amount of errors in all your answers.

The following data may be used:
\begin{itemize}
\item $d_{\mathrm{Sirius,Earth}} = 2.6$\,pc --- the Sirius--Earth distance;
\item $m_{\mathrm{Sirius,Earth}} = -1.46^{\mathrm{m}}$ --- the apparent
  magnitude of Sirius measured from the Earth;
\item $d_{\mathrm{Sun,Earth}} = 1$\,AU --- the Sun--Earth distance;
\item $m_{\mathrm{Sun,Earth}} = -26.78^{\mathrm{m}}$ --- the apparent
  magnitude of the Sun as seen from Earth;
\item $d_{\mathrm{Sirius,Planet}} = 10$\,AU --- distance between Sirius and
  its planet where the observer A is located.
\end{itemize}

In the table below information for the stars from the Mizar system as
measured from the Earth is given.

\begin{center}
\begin{tabular}{cccc}
\hline
Star & Name of the & Apparent & Parallax \\
number & star & magnitude & (milli arc seconds) \\
\hline
1 & Alcor & $3.99 \pm 0.01$ & $39.91 \pm 0.13$ \\
2 & Mizar A & $2.23 \pm 0.01$ & $38.01 \pm 1.71$ \\
3 & Mizar B & $3.86 \pm 0.01$ & $38.01 \pm 1.71$ \\
4 & Sidus Ludoviciana & $7.56 \pm 0.01$ & $8 \pm 4$ \\
\hline
\end{tabular}
\end{center}

The equatorial coordinates of Mizar system ($\sigma_1$) and respectively of
Sirius ($\sigma_2$), located on the heliocentric map are:
\[ \alpha_{\mathrm{Mizar}} = \alpha_1 =
   13^{\mathrm{h}}23^{\mathrm{min}}55.5^{\mathrm{s}}; \quad
   \delta_{\mathrm{Mizar}} = \delta_1 = 54^\circ 55' 31''; \quad
   \alpha_{\mathrm{Sirius}} = \alpha_2 = 6^{\mathrm{h}}45^{\mathrm{min}};
   \quad
   \delta_{\mathrm{Sirius}} = \delta_2 = -16^\circ 43'. \]

Note
\[ \ln(1-x) \approx -x \text{ for } x \ll 1 \]
\[ e^x \approx 1 + x \text{ for } x \ll 1 \]
\ioaafig{2014_DA_D3-f1.pdf}

\probhead{6.30}{Interstellar extinction / IC 342}{2015}{Magelang, Indonesia}{DA D2}{}
% source id: 2015_DA_D2   tag: extinction curve, Cepheid distance IC 342
BVRIJHKLMN photometry of 2 stars from the constellation Cassiopeia is given in
Table 5. For both stars it is believed that their light is affected by
extinction by diffuse Interstellar Medium (ISM) only. Assuming that the
observation is done from outside the atmosphere.

\begin{description}
\item[a)] Using the data given in Tables 5 to 9, plot $E_{X-V}/E_{B-V}$ as a
  function of $1/\lambda_X$ for filters B, V, R, I, J, H, K, L, M, N for both
  stars. Fit approximate curves by eye (in particular, note that
  $E_{X-V}/E_{B-V} \sim \mathit{const.}$ as $1/\lambda_X \to 0$). X is each
  band in the photometric system.

  $E_{B-V}$ is the colour excess.
\item[b)] Using the graphs obtained in a), estimate $R_V$ and $R_R$ for each
  star.
  \[ R_V = \frac{A_V}{E_{B-V}} \quad \text{and} \quad
     R_R = \frac{A_R}{E_{R-I}} \]
  ($A_V$ is the absorption in V).
\end{description}

Now apply these results in order to derive a distance estimate for IC 342, a
spiral galaxy in Cassiopeia obscured by Milky Way. You should assume that the
properties of the ISM in IC 342 are similar to those of the ISM in our Galaxy.

\begin{description}
\item[c)] Using the period-magnitude diagrams for 20 Cepheids from IC 342
  (Figures 2 and 3) and assuming the period-luminosity relations:
  \[ \langle M_R \rangle = -2.91 \left( \log\left(\frac{P}{\mathrm{day}}\right)
     - 1 \right) - 4.04 \quad \text{and} \quad
     \langle M_I \rangle = -3.00 \left( \log\left(\frac{P}{\mathrm{day}}\right)
     - 1 \right) - 4.06 \]
  where $\langle M_R \rangle$ and $\langle M_I \rangle$ are the mean absolute
  magnitudes in filters R and I, find $A_R$ for objects in IC 342. Find the
  distance to IC 342.
\end{description}

\medskip
\noindent Table 5: BVRIJHKLMN photometry of two stars in Cassiopeia

\begin{center}
\begin{tabular}{lccccccccccc}
\hline
Star & MK & $B$ & $V$ & $R$ & $I$ & $J$ & $H$ & $K$ & $L$ & $M$ & $N$ \\
 & class & mag & mag & mag & mag & mag & mag & mag & mag & mag & mag \\
\hline
HD 4817 & K3Iab & 8.08 & 6.18 & 4.73 & 3.64 & 2.76 & 1.86 & 1.54 & 1.32 &
  1.59 & -- \\
HD 11092 & K4II & 8.66 & 6.57 & -- & -- & 3.10 & 2.14 & 1.63 & 1.41 & 1.65 &
  1.44 \\
\hline
\end{tabular}
\end{center}

\medskip
\noindent Table 6: $(B-V)_0$ intrinsic colours for selected sp.\ types and
luminosity classes

\begin{center}
\begin{tabular}{ccc}
\hline
 & \multicolumn{2}{c}{$(B-V)_0$ [mag]} \\
 & II & Iab / Ia \\
\hline
F0 & -- & 0.15 \\
G0 & 0.73 & 0.82 \\
K0 & 1.06 & 1.18 \\
K3 & 1.40 & 1.42 \\
K4 & 1.42 & 1.50 \\
\hline
\end{tabular}
\end{center}

\medskip
\noindent Table 7: Infrared intrinsic colours for selected sp.\ types of
supergiant stars

\begin{center}
\begin{tabular}{ccccccccc}
\hline
 & $(V-R)_0$ & $(V-I)_0$ & $(V-J)_0$ & $(V-H)_0$ & $(V-K)_0$ & $(V-L)_0$ &
   $(V-M)_0$ & $(V-N)_0$ \\
 & mag & mag & mag & mag & mag & mag & mag & mag \\
\hline
F0 & 0.20 & 0.31 & 0.36 & 0.51 & 0.60 & 0.64 & 0.65 & 0.82 \\
G0 & 0.55 & 0.90 & 1.14 & 1.52 & 1.71 & 1.72 & 1.72 & 1.98 \\
K0 & 0.95 & 1.59 & 2.01 & 2.64 & 2.80 & 2.87 & 2.79 & 3.14 \\
K3 & 1.13 & 1.96 & 2.41 & 3.14 & 3.25 & 3.39 & 3.25 & 3.63 \\
K4 & 1.20 & 2.13 & 2.59 & 3.37 & 3.44 & 3.62 & 3.46 & 3.84 \\
\hline
\end{tabular}
\end{center}

\medskip
\noindent Table 8: Infrared intrinsic colours for selected sp.\ types of
giant stars

\begin{center}
\begin{tabular}{ccccccccc}
\hline
 & $(V-R)_0$ & $(V-I)_0$ & $(V-J)_0$ & $(V-H)_0$ & $(V-K)_0$ & $(V-L)_0$ &
   $(V-M)_0$ & $(V-N)_0$ \\
 & mag & mag & mag & mag & mag & mag & mag & mag \\
\hline
K0 & 0.60 & 1.03 & 1.23 & 1.72 & 1.94 & 1.97 & 1.90 & 1.92 \\
K3 & 0.86 & 1.39 & 1.84 & 2.40 & 2.69 & 2.82 & 2.70 & 2.73 \\
K4 & 0.96 & 1.61 & 2.16 & 2.77 & 3.05 & 3.22 & 3.08 & 3.02 \\
\hline
\end{tabular}
\end{center}

\medskip
\noindent Table 9: Effective wavelengths of selected photometric filters

\begin{center}
\begin{tabular}{ccccccccccc}
\hline
Filter & B & V & R & I & J & H & K & L & M & N \\
$\lambda_{\mathrm{F}}$/nm & 450 & 555 & 670 & 870 & 1200 & 1620 & 2200 &
  3500 & 5000 & 9000 \\
\hline
\end{tabular}
\end{center}

% Fig. 2: <R> is the mean apparent magnitude in filter R
\ioaafig[0.5]{2015_DA_D2-f5.pdf}
% Fig. 3: <I> is the mean apparent magnitude in filter I
\ioaafig[0.5]{2015_DA_D2-f6.pdf}

\section*{Observation --- written and sky map}

\probhead{6.31}{Naked-eye magnitude estimates near Cepheus}{2009}{Tehran, Iran}{OBS-SKY 2}{80}
% source id: 2009_OBS-SKY_2   tag: naked-eye magnitude comparison
\begin{description}
\item[2.1:] Figure 2 shows a part of the sky which contains \textbf{Cephei
  constellation}, for 22 October 2009 at 22:00 local time. Five bright stars
  in Cephei constellation are identified by numbers (1, 2, \dots, 5) and
  common names. Estimate the angular distances (in units of degrees) between
  two pairs of stars shown in table 2-1 and complete this table with your
  answers. \textbf{(40 Points)}

  \medskip
  \noindent Table 2-1: Angular Distance

  \begin{center}
  \begin{tabular}{lc}
  \hline
  Pairs of stars & Angular Distance (degrees) \\
  \hline
  1 (Errai) and 2 (Alfirk) & \\
  1 (Errai) and 3 (Alderamin) & \\
  \hline
  \end{tabular}
  \end{center}

\item[2.2:] Use table 2-2 and figure 2, then Estimate the ``apparent visual
  magnitude'' of stars 2 (Alfirak) and 3 (Alderamin) and complete Table 2-3.
  % sic: "Alfirak" as printed in the source; the star is called Alfirk
  % elsewhere in the same problem
  \textbf{(40 Points)}

  \medskip
  \noindent Table 2-2

  \begin{center}
  \begin{tabular}{cc}
  \hline
  Star Name & Apparent Visual Magnitude \\
  \hline
  Polaris & 1.95 \\
  Altais & 3.05 \\
  Segin & 3.34 \\
  \hline
  \multicolumn{2}{c}{All of these stars, are marked in the figure 2} \\
  \hline
  \end{tabular}
  \end{center}

  \medskip
  \noindent Table 2-3: Magnitude Estimation

  \begin{center}
  \begin{tabular}{ccc}
  \hline
  Star Number & Star Name & Apparent Visual Magnitude \\
  \hline
  2 & Alfirk & \\
  3 & Alderamin & \\
  \hline
  \end{tabular}
  \end{center}
\end{description}

\ioaafig{2009_OBS-SKY_2-f1.pdf}
% ($IOAA_PAPERS_DIR/observation/2009_OBS.pdf, page 7): sky chart around
% Cepheus with stars 1 (Errai), 2 (Alfirk), 3 (Alderamin), 5 (Iota Cephei),
% plus Polaris, Altais and Segin marked.

\section*{Observation --- telescope}

\probhead{6.32}{Magnitude of a missing star in NGC 6231}{2012}{Rio de Janeiro, Brazil}{OBS-TEL 2}{}
% source id: 2012_OBS-TEL_2   tag: visual magnitude estimation
Use \emph{star chart-2} to estimate the magnitude of the missing star, shown
as a cross, inside NGC 6231. Use the magnitude of other stars as reference.

\textbf{Note}: To avoid confusion between decimal dots and real stars, dots
where supressed. % sic: "where supressed" as printed in the source
So, magnitude 60 corresponds to magnitude 6.0. Give your
answer using one decimal figure and 0.1 precision.

\ioaafig{2012_OBS-TEL_2-f1.pdf}
% ($IOAA_PAPERS_DIR/observation/2012_OBS.pdf, page 3): star chart of
% NGC 6231 with the missing star marked as a cross and magnitudes of
% reference stars printed without decimal points.

\clearpage
\section*{Solutions to Chapter 6}

\solhead{6.1}{}
% source id: 2007_TH_1.10
\begin{align*}
\text{flux density} &= \frac{\text{luminosity}}{4\pi \times (\text{distance})^2}\\
\text{distance} &= \sqrt{\frac{0.4 \times 3.826 \times 10^{26}}{4\pi \times 6.23 \times 10^{-14}}}
  = 1.398 \times 10^{19}~\mathrm{m}\\
&= \frac{1.398 \times 10^{19}}{3.0856 \times 10^{16}}~\mathrm{pc} = 453~\mathrm{pc}
\end{align*}

\solhead{6.2}{}
% source id: 2007_TH_1.11
\[
L_{\text{supernova}} = 10^{10} L_{\odot}
\]
Let $D$ be the distance to that supernova.

Then the flux intensity on Earth would be
$\displaystyle I = \frac{10^{10} L_{\odot}}{4\pi D^2}$.

This must be the same as $\displaystyle \frac{L_{\odot}}{4\pi \left(D_{\oplus-\odot}\right)^2}$
\[
\frac{10^{10} L_{\odot}}{4\pi D^2} = \frac{L_{\odot}}{4\pi\,(1~\mathrm{A.U.})^2}
\]
\begin{align*}
\therefore\quad D &= 10^{\frac{10}{2}}~\mathrm{A.U.} = 10^{5}~\mathrm{A.U.}\\
&= \frac{1.496 \times 10^{16}}{3.0856 \times 10^{16}}~\mathrm{pc} = 0.485~\mathrm{pc}\\
&= 0.485 \times 3.2615~\mathrm{ly} = 1.58~\mathrm{ly}
\end{align*}

\solhead{6.3}{}
% source id: 2007_TH_1.12
The difference in frequencies is due to the relativistic Doppler shift. Since
the observed frequency of emission from the gas cloud is higher than the
laboratory frequency $\nu_0$, the gas cloud must be approaching the observer.
Hence,
\begin{align*}
\nu &= \nu_0 \sqrt{\frac{c+v}{c-v}}\\
\frac{v}{c} &= \frac{\left(\dfrac{\nu}{\nu_0}\right)^2 - 1}
                    {\left(\dfrac{\nu}{\nu_0}\right)^2 + 1}\\
&= \frac{0.001752379}{2.001752379} = 0.000875422\\
v &= \left(0.000875422\right)\left(2.99792458 \times 10^{8}~\mathrm{m.s^{-1}}\right)\\
&= 0.00262445 \times 10^{8}~\mathrm{m.s^{-1}}\\
&= 262.445~\mathrm{km.s^{-1}}
\end{align*}

\solhead{6.4}{}
% source id: 2007_TH_1.15
The flux ratio versus magnitude difference implies
\begin{align*}
m_1 - m_2 &= -2.5 \log\!\left(\frac{f_1}{f_2}\right)\\
\frac{f_1}{f_2} &= 10^{\frac{(m_2 - m_1)}{2.5}}
\end{align*}
So for a magnitude difference of $(-1.5) - 6 = -7.5$ we find a flux ratio of
\[
\frac{f_{\min}}{f_{\max}} = 10^{\frac{-7.5}{2.5}} = 10^{-3}
\]
And thus the visible stars range only over a factor of 1000 in brightness.

\solhead{6.5}{}
% source id: 2007_TH_1.8
From the Wien's displacement law
\begin{align*}
\lambda_{\max} T &= 2.8977 \times 10^{-3}~\mathrm{m.K}\\
\lambda_{\max} &= \frac{2.8977 \times 10^{-3}}{4000}~\mathrm{m}
  = 7.244 \times 10^{-7}~\mathrm{m}\\
&= 724~\text{nanometers}
\end{align*}

\solhead{6.6}{}
% source id: 2008_TH_L2
$U = 8.15$, $B = 8.50$, and $V = 8.14$.\\
$(U-B)_{\mathrm{o}} = -0.45$. \quad
$R = 2.3\,R_\odot = 1.60 \times 10^{11}$\,cm,\\
$M_{bol} = -0.25$, $BC = -0.15$

\begin{description}
\item[a)] $U - B = 8.15 - 8.50 = -0.35$

The color excess for $U-B$:
\[
E(U-B) = (U-B) - (U-B)_{\mathrm{o}} = -0.35 - (-0.45) = 0.10
\]
The relation between color excess of $U-B$ and of $B-V$:
\[
E(U-B) = 0.72\,E(B-V)
\;\Longrightarrow\;
E(B-V) = 0.10/0.72 = 0.14 \tag*{(15\%)}
\]
Let $A_v$ be the interstellar extinction and $R = 3.2$, then
\[
A_v = 3.2\,E(B-V) = 3.2\,(0.14) = 0.45
\]
\[
V - V_{\mathrm{o}} = A_v
\;\Longrightarrow\;
V_{\mathrm{o}} = V - A_v = 8.14 - 0.45 = 7.69
\]
\begin{align*}
E(B-V) = (B-V) - (B-V)_{\mathrm{o}}
\;\Longrightarrow\;
(B-V)_{\mathrm{o}} &= (B-V) - E(B-V)\\
&= (8.50 - 8.14) - 0.14 = 0.22 \tag*{(15\%)}
\end{align*}
\[
(B-V)_{\mathrm{o}} = B_{\mathrm{o}} - V_{\mathrm{o}} = 0.22
\;\Longrightarrow\;
B_{\mathrm{o}} = 0.22 + V_{\mathrm{o}} = 0.22 + 7.69 = 7.91 \tag*{(10\%)}
\]
\[
(U-B)_{\mathrm{o}} = U_{\mathrm{o}} - B_{\mathrm{o}} = -0.45
\;\Longrightarrow\;
U_{\mathrm{o}} = B_{\mathrm{o}} - 0.45 = 7.91 - 0.45 = 7.46 \tag*{(10\%)}
\]

\item[b)] Luminosity -- absolute magnitude relation:
\begin{align*}
M_{bol} - M_{bol\odot} = -2.5 \log L/L_\odot
\;\Longrightarrow\;
L/L_\odot &= 10^{-(M_{bol} - M_{bol\odot})/2.5}\\
&= 10^{-(-0.25 - 4.75)/2.5}\\
&= 100 \tag*{(10\%)}
\end{align*}
\[
L = 100\,L_\odot = 3.90 \times 10^{35}~\mathrm{erg/det}
% sic: "det" (detik, Indonesian for second) in the source
\]
\[
L = 4\pi\sigma R^2\,T_{ef}^4
\]
\[
T_{ef} = \left( \frac{L}{4\pi\sigma R^2} \right)^{\!1/4}
  = \left( \frac{3.90 \times 10^{35}}
                {4\pi (5.67 \times 10^{-5})(1.60 \times 10^{11})^2}
    \right)^{\!1/4}
  = 12\,092~\mathrm{K} \tag*{(15\%)}
\]

\item[c)]
\[
M_v - M_{bol} = BC
\;\Longrightarrow\;
M_v = M_{bol} + BC = -0.25 - 0.15 = -0.40 \tag*{(10\%)}
\]
\begin{align*}
m_v - M_v = -5 + 5 \log d + A_v
\;\Longrightarrow\;
5 \log d &= m_v - M_v + 5 - A_v\\
&= 8.14 + 0.40 + 5 - 0.45 = 13.09
\end{align*}
\[
d = 10^{13.09/5} = 10^{2.62} = 414.9~\mathrm{pc} \tag*{(15\%)}
\]
\end{description}

\solhead{6.7}{}
% source id: 2009_TH_P2
To calculate flux of a $m = 6$ star we use the Sun as standard candle
\[
m_1 - m_2 = -2.5 \log\frac{f_1}{f_2} \tag*{(1 point)}
\]
\[
6 - (-26.8) = -2.5 \log\frac{f_1}{1.37 \times 10^{3}}
\]
\[
f_1 = 1.04 \times 10^{-10}~(\mathrm{w/m^2}) \tag*{(3 points)}
\]
We need to know how much energy a visual photon has (at 550\,nm)
\begin{align*}
E_p = h\nu = \frac{hc}{\lambda}
&= \frac{6.63 \times 10^{-34} \times 3.00 \times 10^{8}}{550 \times 10^{-9}}\\
&= 3.62 \times 10^{-19}~J \tag*{(3 points)}
\end{align*}
Then number of photon which arrive to our eye per second is
\[
N = \frac{f_1 \pi r_e^2}{E_p}
  = \frac{1.04 \times 10^{-10}}{3.62 \times 10^{-19}} \times 3.1416 \times 0.003^2
  = 8 \times 10^{3}~s^{-1} \tag*{(3 points)}
\]

\solhead{6.8}{}
% source id: 2010_TH_S1
Let $F_1$, $F_2$, and $F_0$ be the flux of the first, the second and the binary
system, respectively.
\begin{align*}
\Delta m &= -2.5 \lg (F_1/F_2)\\
(1-2) &= -2.5 \lg (F_1/F_2) \tag*{5}
\end{align*}
So, $F_1/F_2 = 10^{1/2.5} = 10^{0.4}$
\begin{equation*}
F_0 = F_1 + F_2 = F_1 (1 + 10^{-0.4}) \tag*{3}
\end{equation*}
The magnitude of the binary $m$ is:
\begin{equation*}
m - 1 = -2.5 \lg (F_0/F_1) = -2.5 \lg (F_1 (1+0.398)/F_1) = -0.36^m \tag*{2}
\end{equation*}
So, $m = 0.64^m$

\solhead{6.9}{}
% source id: 2011_TH_S8
Gravitational force can be expressed in the form:
$F_g = \dfrac{4}{3}\pi \cdot r^3 \cdot \rho \cdot G \cdot \dfrac{M}{R^2}$ \hfill[2]

The force of radiation reads $F_r = \pi \cdot r^2 \cdot C_0/c$ \hfill[2]\\
where $r, \rho, R, G, M, C_0$ are a radius and density of dust, distance form % sic: "form" for "from"
the Sun, gravitational constant, Solar mass and Solar constant respectively.

Comparing these forces one obtains $F_g = F_r$ \hfill[2]

and finally one gets the dust grain radius
$r = \dfrac{3}{4}\,\dfrac{C_0 \cdot R^2}{c \cdot \rho \cdot G \cdot M}$ \hfill[2]

Substituting numerical values one obtains $r \approx 5 \cdot 10^{-7}$\,m \hfill[2]

\solhead{6.10}{}
% source id: 2011_TH_S9
Finding that surface brightness is independent of distance \hfill[2]

Relationship between $m$ and the Sun surface brightness:
$m - \mu_\odot = -2.5 \log\left( (\pi\cdot\theta^2/4) \,/\, 1~\mathrm{arcsec}^2 \right)$. \hfill[2]

Calculations of $\mu_\odot = -10.79$ \hfill[1]

Galaxy vs.\ Sun surface brightness (in magnitudes$\cdot$arcsec$^{-2}$)\\
$\mu - \mu_\odot = 18.0 + 10.79 = 28.79$ \hfill[1]

Relationship between the sought sky fraction covered by stars, $x$,
and the surface brightness ratio: $\mu - \mu_\odot = -2.5 \log x$ \hfill[3]

Calculations $x = 10^{-0.4(\mu - \mu_\odot)} \approx 3.1 \cdot 10^{-12}$ \hfill[1]

\solhead{6.11}{}
% source id: 2012_TH_L1
\begin{description}
\item[(1.1)] Look at the figure:

\ioaafig{2012_TH_L1-s1.pdf}
\[
\tan\alpha \approx \alpha \approx \sin\alpha \approx \frac{2R}{D}
\]
\[
\Delta m = m_B - m_A = -2.5 \log\frac{F_B}{F_A}
\]
The brightness as seen by the astronomer ($F_A$) is the brightness of $N$
stars with absolute magnitude $M_0$ (and luminosity $L_0$) at a distance $D$
from him:
\[
F_A = N \frac{L_0}{4\pi D^2} \tag*{(4 pt)}
\]
We assume the cluster is homogeneous, so the density of stars $\rho$ inside
the cluster should be constant. The brightness coming from a spherical shell
with small thickness $\Delta R$ at a distance $R'$ is
\[
F_{\text{shell}}(R', \Delta R)
  = \rho V_{\text{shell}} \frac{L_0}{4\pi R'^2}
  = \frac{N}{\frac{4}{3}\pi R^3}\, 4\pi R'^2 \Delta R\, \frac{L_0}{4\pi R'^2}
  = \frac{3 N L_0 \Delta R}{4\pi R^3} \tag*{(5 pt)}
\]
The brightness doesn't depend on the radius of the shell $R'$. If we divide
the cluster in $n$ shells, $R = n \Delta R$, the brightness of the cluster as
seen by the astronomer 2 is: % sic: "astronomer 2" in the source; this is the
% flux seen by the observer at the centre (the biologist)
\[
F_B = n F_{\text{shell}} = n \frac{3 N L_0 \Delta R}{4\pi R^3}
    = \frac{3 N L_0}{4\pi R^2} \tag*{(4 pt)}
\]
So the difference in the magnitudes $\Delta m$ is:
\[
\Delta m = -2.5 \log\!\left(\frac{\;\frac{3NL_0}{4\pi R^2}\;}
                                 {\frac{NL_0}{4\pi D^2}}\right)
  = -2.5 \log\!\left[ 3 \left(\frac{D}{R}\right)^{\!2} \right]
  = 2.5 \log\!\left(\frac{\alpha^2}{12}\right) \tag*{(2 pt)}
\]

\item[(1.2)] The total energy received is proportional to the objective area:
\[
F'_A = F_A \frac{\pi \left(\frac{x}{2}\right)^2}
                {\pi \left(\frac{d_{\mathrm{pup}}}{2}\right)^2}
\quad\Rightarrow\quad
x = 6 \cdot 10^{-3} \sqrt{\frac{F'_A}{F_A}} \tag*{(3 pt)}
\]
\[
F'_A = F_B \quad\Rightarrow\quad
\frac{F'_A}{F_A} = \frac{F_B}{F_A} = \frac{12}{\alpha^2}
\quad\Rightarrow\quad
x = 6 \cdot 10^{-3} \sqrt{\frac{12}{\alpha^2}}
  = \frac{0.021}{\alpha}~m \tag*{(3 pt)}
\]

\item[(1.3)] In this case, a fraction
$\xi(d) = \dfrac{A(\alpha)}{4\pi d^2}$ of the light that comes from the
distance $d$ is seen by the biologist. $A(\alpha)$ is the area of the
interior of a cone with aperture $\alpha$ that intersects the sphere with
radius $d$ and center in the vertices of the cone (a spherical cap of radius
$d$ and height $D$): \hfill(2 pt)
\[
\cos\left(\frac{\alpha}{2}\right) = \frac{d - D}{d}
\quad\text{and}\quad
A(\alpha) = 2\pi D d
\;\Rightarrow\;
A(\alpha) = 2\pi d^2 \left[ 1 - \cos\left(\frac{\alpha}{2}\right) \right]
\;\Rightarrow\;
\xi = \sin^2\left(\frac{\alpha}{4}\right) \tag*{(4 pt)}
\]
So the fraction is independent of the distance $d$, and the fraction of light
that the biologist sees is $F'_B = \xi F_B$:
\[
\Delta m' = m'_B - m_A = -2.5 \log\frac{F'_B}{F_A}
  = -2.5 \log\!\left[ \frac{12 \sin^2\left(\frac{\alpha}{4}\right)}{\alpha^2} \right]
\tag*{(3 pt)}
\]
\end{description}

\solhead{6.12}{}
% source id: 2012_TH_S9
The intrinsic color indexes are $U - V = 0.300$
% sic: the source says "U - V = 0.300"; the problem gives B - V = 0.300
and $U - V = 0.330$, therefore
\[
U - B = (U - V) - (B - V) = 0.030
\]
The difference between the real and the measured color indexes are due to the
absorption in the interstellar medium, as well as due the stronger absorption
of the planetary nebula (in fact, during the calculations, one can easily see
that the interstellar absorption is not enough to account for the measured
values). So, for the first (nearest) star:
\begin{align*}
(B-V)_1 &= (B-V) + (E_B - E_V)\cdot D + (E'_B - E'_V)\cdot R = 0.327\\
(U-V)_1 &= (U-V) + (E_U - E_V)\cdot D + (E'_U - E'_V)\cdot R = 0.038\\
% sic: the values 0.038 and 0.365 on these two lines are swapped in the
% source; with the given data (U-V)_1 = 0.365 and (U-B)_1 = 0.038
(U-B)_1 &= (U-B) + (E_U - E_B)\cdot D + (E'_U - E'_B)\cdot R = 0.365
\tag*{(3 pt)}
\end{align*}
(the terms being, in order: apparent, intrinsic, interstellar absorption, and
nebula absorption). Calculating the interstellar absorption, we can find for,
the nebula values:
\begin{align*}
(E'_B - E'_V)\cdot R &= 0.0155\\
(E'_U - E'_V)\cdot R &= 0.0100\\
(E'_U - E'_B)\cdot R &= -0.0055 \tag*{(1 pt)}
\end{align*}
So, for the second (farthest) star,
\begin{align*}
(B-V)_2 &= (B-V) + (E_B - E_V)\cdot 3D + (E'_B - E'_V)\cdot 2R = \mathbf{0.3655}\\
(U-V)_2 &= (U-V) + (E_U - E_V)\cdot 3D + (E'_U - E'_V)\cdot 2R = \mathbf{0.425}\\
(U-B)_2 &= (U-B) + (E_U - E_B)\cdot 3D + (E'_U - E'_B)\cdot 2R = \mathbf{0.0595}
\tag*{(4 pt + 2 pt final answers)}
\end{align*}

\solhead{6.13}{}
% source id: 2013_TH_S13
The absorption due to the fog in London is obviously
\[
A = -26.8 - (-12.74) = -14.06~\text{mag}. \tag*{(2 Point)}
\]
Rearranging the equation $I_\nu(r) = I_\nu(0) \times e^{-\tau}$, we get
\[
\frac{I_\nu(0)}{I_\nu(r)} = e^{\tau}, \tag*{(5 Points)}
\]
or
\[
A = \Delta m = 2.5 \times \log(e) \times (-\tau)
\quad\text{or}\quad
\tau = (-A)/(2.5 \times \log e) = 14.06/1.08
\quad\text{or}\quad % sic: source has the Greek word for "or" here
\tau = 12.9 \tag*{(3 Points)}
\]

\solhead{6.14}{}
% source id: 2013_TH_S8
The apparent magnitude of star (a) is $m_a$, of star (b) is $m_b$ and that of
the system as a whole is $m_{a+b}$. The corresponding apparent brightnesses are
$\ell_a$, $\ell_b$ and $\ell_{a+b} = \ell_a + \ell_b$.

For star (a):
\[
m_{a+b} - m_a = -2.5 \log\!\left(\frac{\ell_a + \ell_b}{\ell_a}\right)
\]
and because $\dfrac{\ell_b}{\ell_a} = \dfrac{1}{2}$, we get
$m_a = m_{a+b} + 2.5 \log\left(1 + \tfrac{1}{2}\right)$ or
$m_a = 5 + 2.5 \log (3/2)$ and finally $m_a = 5.44$~mag. \hfill(5 point)

Similarly for star (b):
\[
m_{a+b} - m_b = -2.5 \log\!\left(\frac{\ell_a + \ell_b}{\ell_b}\right)
\]
and because $\dfrac{\ell_a}{\ell_b} = 2$, we get
$m_b = m_{a+b} + 2.5 \log (3)$ or
$m_b = 5 + 2.5 \log (3)$ and finally $m_b = 6.19$~mag. \hfill(5 Point)

\solhead{6.15}{Space-ship orbiting the Sun}
% source id: 2014_TH_P11
According to the Stefan--Boltzmann law, the luminosity of the Sun is:
\begin{equation*}
L_{\text{sun}} = Q_{\text{sun}} \cdot 4\pi R_{\text{Sun}}^2
  = \sigma T_{\text{Sun}}^4 \cdot 4\pi R_{\text{Sun}}^2, \tag{1}
\end{equation*}
At distance $d$ from the Sun, where the space ship is the energy which passes
the unit of surface in an unit of time is:
\begin{equation*}
\phi_{\text{Sun,d}} = \frac{L_S}{4\pi d^2}
  = \frac{\sigma T_S^4 \cdot 4\pi R_S^2}{4\pi d^2}. \tag{2}
\end{equation*}
The space ship receive through its entire surface, in the unit of time, the
energy:
\[
P_{\text{received}} = \frac{\sigma T_{\text{Sun}}^4 \cdot 4\pi R_{\text{Sun}}^2}{4\pi d^2}
  \cdot \pi R_{\text{ship}}^2.
\]
Corresponding to its temperature, $T_N$, according the Stefan--Boltzmann law,
the emitted energy by starship through its hole surface % sic: "hole" for "whole"
in the unit of time:
\[
P_{\text{emis,N}} = \sigma T_N^4 \cdot 4\pi R_N^2.
\]
When the temperature stabilized at thermic equilibrium:
\begin{align*}
P_{\text{received,N}} &= P_{\text{emis,N}}\\
\frac{\sigma T_S^4 \cdot 4\pi R_S^2}{4\pi d^2} \cdot \pi R_N^2
  &= \sigma T_N^4 \cdot 4\pi R_N^2
\end{align*}
the distance of orbiting the Sun of the space ship is:
\[
d = \frac{T_S^2 R_S}{2 T_N^2},
\]
The angular diameter of the Sun as seen from the space ship:
\begin{align*}
\alpha/2 &= \frac{R_S}{d}\\
\alpha &= \frac{2 R_S}{d} = 4 \left( \frac{T_N}{T_S} \right)^{\!2}
\end{align*}
According the Pogson formula written for Sun seen from Earth and space ship
the following relation occurs:
\begin{equation*}
\lg \frac{E_{\text{S,Nava}}}{E_{\text{S,P}}} = -0.4\,(m - m_0) \tag{2}
\end{equation*}
% sic: the source labels this equation (2) although (2) was already used above
\[
2 \cdot \lg\!\left( \frac{d_0}{d} \right) = -0.4\,(m - m_0)
\]
The apparent magnitude of the Sun as seen from the space ship
\[
m = m_0 - 5^{\mathrm{m}} \cdot \lg \frac{2 d_0 T_N^2}{R_S T_S^2}
\]

\solhead{6.16}{The effective temperature on the surface of a star}
% source id: 2014_TH_P7
\begin{description}
\item[a.] We start from the definition of $r$:
\[
r = \frac{\Delta E}{\Delta t \cdot S_{\text{stea}} \cdot \Delta\lambda}
  = \frac{\Delta E}{\Delta t \cdot 4\pi R^2 \cdot \Delta\lambda}
\quad\text{where $R$ is the radius of the star}
\]
% "S_stea" (Romanian: star surface) as in the source
\[
r = \frac{2\pi h c^2}{\lambda^5 \left( e^{hc/k\lambda T} - 1 \right)}
\]
Considering $d$ as the distance from the star to the Earth, the
definition-relation of the spectral intensity can be written as follows:
\begin{align*}
I(\lambda) &= \frac{\;\dfrac{\Delta E}{\Delta t \cdot \Delta\lambda}\;}{4\pi d^2}\\
I(\lambda) &= \frac{2\pi h c^2 R^2}
                  {d^2 \lambda^5 \left( e^{hc/\lambda k T} - 1 \right)};
\end{align*}
Particularly for each wavelength:
\[
I_1(\lambda_1) = \frac{2\pi h c^2 R^2}
  {d^2 \lambda_1^5 \left( e^{hc/\lambda_1 k T} - 1 \right)};\qquad
I_2(\lambda_2) = \frac{2\pi h c^2 R^2}
  {d^2 \lambda_2^5 \left( e^{hc/\lambda_2 k T} - 1 \right)};
\]
The ratio of the 2 above relations
\begin{equation*}
\frac{I_1(\lambda_1)}{I_2(\lambda_2)}
  = \left( \frac{\lambda_2}{\lambda_1} \right)^{\!5}
    \cdot \frac{e^{hc/\lambda_2 k T} - 1}{e^{hc/\lambda_1 k T} - 1}
\tag{1}
\end{equation*}
Represents an equation which allow to find out the temperature of star's
surface $T$ by using spectral measurements.

\item[b.] If we consider that $hc \gg \lambda k T$, then:
\[
e^{hc/\lambda_1 k T} - 1 \approx e^{hc/\lambda_1 k T}
\quad\text{and}\quad
e^{hc/\lambda_2 k T} - 1 \approx e^{hc/\lambda_2 k T}
\]
The relation (1) becomes
\[
\frac{I_1(\lambda_1)}{I_2(\lambda_2)}
  = \left( \frac{\lambda_2}{\lambda_1} \right)^{\!5}
    \cdot \frac{e^{hc/\lambda_2 k T}}{e^{hc/\lambda_1 k T}}
  = \left( \frac{\lambda_2}{\lambda_1} \right)^{\!5}
    \cdot e^{hc/\lambda_2 k T - hc/\lambda_1 k T}
\]
\[
\frac{I_1(\lambda_1)}{I_2(\lambda_2)}
  \cdot \left( \frac{\lambda_1}{\lambda_2} \right)^{\!5}
  = e^{hc/\lambda_2 k T - hc/\lambda_1 k T}
\]
\[
T = \frac{hc\,(\lambda_1 - \lambda_2)}
         {k \lambda_1 \lambda_2 \cdot
          \ln\!\left[ \dfrac{I_1(\lambda_1)}{I_2(\lambda_2)}
          \cdot \left( \dfrac{\lambda_1}{\lambda_2} \right)^{\!5} \right]}
\]

\item[c.] If $hc \ll \lambda k T$, then:
\begin{align*}
\frac{hc}{k\lambda_1 T} \ll 1;\quad
e^{hc/\lambda_1 k T} - 1 &\approx 1 + \frac{hc}{k\lambda_1 T} - 1
  = \frac{hc}{k\lambda_1 T}\\
\frac{hc}{k\lambda_2 T} \ll 1;\quad
e^{hc/\lambda_2 k T} - 1 &\approx 1 + \frac{hc}{k\lambda_2 T} - 1
  = \frac{hc}{k\lambda_2 T}
\end{align*}
The relation (1) becomes:
\[
\frac{I_1(\lambda_1)}{I_2(\lambda_2)}
  = \left( \frac{\lambda_2}{\lambda_1} \right)^{\!5}
    \cdot \frac{\;\dfrac{hc}{k\lambda_2 T}\;}{\dfrac{hc}{k\lambda_2 T}}
% sic: the source prints k lambda_2 T in both numerator and denominator;
% the denominator should be k lambda_1 T
  = \left( \frac{\lambda_2}{\lambda_1} \right)^{\!4}
\]
\[
I_1 = 2 I_2;\quad
\left( \frac{\lambda_2}{\lambda_1} \right)^{\!4} = 2;\quad
\lambda_2 = \lambda_1 \cdot \sqrt[4]{2} \approx 1.2 \cdot \lambda_1.
\]
\end{description}

\solhead{6.17}{Gradient temperatures}
% source id: 2014_TH_P8
\textbf{Detailed solution}

By using the Plank expression the spectral intensities for the two wavelengths
in the doublet emitted by the star number 1 are:
\[ r_1(\lambda_1) = \frac{2\pi h c^2}
     {\lambda_1^5\left(e^{hc/\lambda_1 k T_1}-1\right)};\quad
   r_1(\lambda_2) = \frac{2\pi h c^2}
     {\lambda_2^5\left(e^{hc/\lambda_2 k T_1}-1\right)}; \]
And by considering $hc \gg k\lambda T$;
\[ r_1(\lambda_1) = \frac{2\pi h c^2}{\lambda_1^5}\,
     e^{-hc/\lambda_1 k T_1};\quad
   r_1(\lambda_2) = \frac{2\pi h c^2}{\lambda_2^5}\,
     e^{-hc/\lambda_2 k T_1}. \tag*{(1)} \]
Respectively, the spectral intensities for the two wavelengths in the doublet
emitted by the star number 2 are:
\[ r_2(\lambda_1) = \frac{2\pi h c^2}
     {\lambda_1^5\left(e^{hc/\lambda_1 k T_2}-1\right)};\quad
   r_2(\lambda_2) = \frac{2\pi h c^2}
     {\lambda_2^5\left(e^{hc/\lambda_2 k T_2}-1\right)}; \]
And by considering $hc \gg k\lambda T$;
\[ r_2(\lambda_1) = \frac{2\pi h c^2}{\lambda_1^5}\,
     e^{-hc/\lambda_1 k T_2};\quad
   r_2(\lambda_2) = \frac{2\pi h c^2}{\lambda_2^5}\,
     e^{-hc/\lambda_2 k T_2}. \tag*{(2)} \]
Using the Pogson formula for star 1 and 2 and for $\lambda_1$, result:
\[ \log\frac{r_1(\lambda_1)}{r_2(\lambda_1)}
   = -0.4\left(m_{1,\lambda_1}-m_{2,\lambda_1}\right)
   = -0.4\cdot\Delta m_{\lambda_1}; \]
\[ \log e^{hc/\lambda_1 k\left(1/T_2-1/T_1\right)}
   = -0.4\cdot\Delta m_{\lambda_1} \tag*{(3)} \]
Similar the Pogson formula for $\lambda_2$:
\[ \log\frac{r_1(\lambda_2)}{r_2(\lambda_2)}
   = -0.4\left(m_{1,\lambda_2}-m_{2,\lambda_2}\right)
   = -0.4\cdot\Delta m_{\lambda_2}; \]
\[ \log e^{hc/\lambda_2 k\left(1/T_2-1/T_1\right)}
   = -0.4\cdot\Delta m_{\lambda_2} \tag*{(4)} \]
Using the relations (3) and (4)
\[ \log e^{hc/\lambda_2 k\left(1/T_2-1/T_1\right)}
   - \log e^{hc/\lambda_1 k\left(1/T_2-1/T_1\right)}
   = -0.4\cdot\Delta m_{\lambda_2} + 0.4\cdot\Delta m_{\lambda_1}; \]
\[ \log e^{hc/\lambda_2 k\left(1/T_2-1/T_1\right)}
   + \log e^{-hc/\lambda_1 k\left(1/T_2-1/T_1\right)}
   = -0.4\cdot\Delta m_{\lambda_2} + 0.4\cdot\Delta m_{\lambda_1}; \]
\[ \left[\frac{hc}{k}\left(\frac{1}{T_2}-\frac{1}{T_1}\right)\cdot
   \left(\frac{1}{\lambda_2}-\frac{1}{\lambda_1}\right)\right]\cdot\log e
   = 0.4\left(\Delta m_{\lambda_1}-\Delta m_{\lambda_2}\right); \]
Because $\log e \approx 0.43$;
$\dfrac{0.4}{\log e}\approx 0.93$;
\[ \frac{1}{T_2}-\frac{1}{T_1}
   = -0.93\cdot\frac{\left(\Delta m_{\lambda_1}-\Delta m_{\lambda_2}\right)}
     {\dfrac{hc}{k}\left(\dfrac{1}{\lambda_2}-\dfrac{1}{\lambda_1}\right)}
   = -0.93\cdot\frac{k\lambda_1\lambda_2
       \left(\Delta m_{\lambda_1}-\Delta m_{\lambda_2}\right)}
     {hc\left(\lambda_1-\lambda_2\right)}; \]
\[ T_1 = T_2\cdot\frac{hc\left(\lambda_1-\lambda_2\right)}
     {hc\left(\lambda_1-\lambda_2\right)
      + 0.93\cdot k\lambda_1\lambda_2 T_2
        \left(\Delta m_{\lambda_1}-\Delta m_{\lambda_2}\right)}. \tag*{(5)} \]

\solhead{6.18}{Pressure of light}
% source id: 2014_TH_P9
The pressure of the emitted radiation is
\[ p_{\mathrm{rad}} = w = \frac{\phi_{planet,D}}{c} \]
\[ p_{\mathrm{rad}}
   = \frac{\sigma T_{planet}^4 R_{planet}^2}{cD^2} \tag*{(1)} \]
% sic: the source writes "planet" subscripts although the problem concerns
% the Sun and a dust particle
As seen in the below image, the pressure due to the solar radiation is
effectively acting on an equivalent plane disc with the diameter $d$ of the
spherical star dust particle
\ioaafig{2014_TH_P9-s1.pdf}

Thus the force acting by the Sun radiation on the star-dust particle is:
\[ F_{\mathrm{rad,S}} = p_{\mathrm{rad,S}}\cdot\pi r^2
   = p_{\mathrm{rad,S}}\cdot\frac{\pi d^2}{4}
   = \frac{\sigma T_{\mathrm{S}}^4 R_{\mathrm{S}}^2}{cL_{\mathrm{S}}^2}
     \cdot\frac{\pi d^2}{4}, \]
% sic: "cL_S^2" in the source; from relation (1) and the following lines the
% denominator should be c D_S^2
The equilibrium condition is
\[ F_{\mathrm{rad,S}} = F_{\mathrm{g}}, \]
Where $F_{\mathrm{g}}$ is the gravitational attraction force between Sun and
the star-dust particle.
\begin{align*}
\frac{\sigma T_{\mathrm{S}}^4 R_{\mathrm{S}}^2}{cD_{\mathrm{S}}^2}
  \cdot\frac{\pi d^2}{4}
  &= K\,\frac{m M_{\mathrm{S}}}{D_{\mathrm{S}}^2};\\
m = \rho V = \rho\,\frac{4\pi r^3}{3}
  &= \rho\,\frac{4\pi}{3}\,\frac{d^3}{8}
   = \rho\,\frac{\pi d^3}{6};\\
\frac{\sigma T_{\mathrm{S}}^4 R_{\mathrm{S}}^2}{cD_{\mathrm{S}}^2}
  \cdot\frac{\pi d^2}{4}
  &= K\rho\,\frac{\pi d^3}{6}\,\frac{M_{\mathrm{S}}}{D_{\mathrm{S}}^2};\\
d &= \frac{3}{2}\cdot\frac{\sigma}{\rho K}\cdot
     \frac{T_{\mathrm{S}}^4 R_{\mathrm{S}}^2}{M_{\mathrm{S}}}.
\end{align*}
% sic: the final formula omits the factor 1/c which follows from the
% preceding line; dimensionally d = (3/2) sigma T_S^4 R_S^2 / (c rho K M_S)

\solhead{6.19}{}
% source id: 2015_TH_S6
Atmospheric optical depth $\tau_A$ can be determined from the measured flux,
that is:
\begin{align*}
S_{\mathrm{meas}} &= \exp(-\tau_A)\, S_{\mathrm{real}}\\
14.27 &= \exp(-\tau_A) \times 21.86
  \;\Rightarrow\; \exp(-\tau_A) = 14.27/21.86 = 0.65
\end{align*}
Then
\[ \tau_A = -\ln(0.65) = 0.43 \]
\hfill(50 points)

Now we have:
\begin{align*}
\tau_A &= \tau_Z \sec z\\
\tau_z &= \tau_A \cos z\\
       &= 0.43 \times \cos(90 - 35)^\circ = 0.25
\end{align*}
where $z$ is zenith angle. \hfill(50 points)

\solhead{6.20}{}
% source id: 2015_TH_S8
During 700\,$\mu$s, the radio wave travels about
$r = ct = 3\times10^{10}\times700\times10^{-6} = 2.1\times10^{7}$\,cm. The
region from where the burst originates must be no larger than the distance
that light can travel during the duration of the burst. So we estimate that
$r$ is the size of the source. ($r = 2.1\times10^{7}$\,cm,
$R = 2.3$\,kpc) \hfill(10 points)

Flux density observed at 1660\,MHz is
\begin{align*}
S_{\mathrm{1660\,MHz}} &= 0.35\ \mathrm{kJy}\\
  &= 0.35\times10^{3}\times10^{-23}\
     \mathrm{ergs\ s^{-1}\ cm^{-2}\ Hz^{-1}}\\
  &= 3.5\times10^{-21}\ \mathrm{ergs\ s^{-1}\ cm^{-2}\ Hz^{-1}}
\end{align*}
\hfill(10 points)

The solid angle subtended by this source of radiation is
\begin{align*}
\Omega = \pi\,(r/R)^2
  &= 3.14\,\bigl(2.1\times10^{7}/2.3\times10^{3}\times
     3.086\times10^{18}\bigr)^2\\
  &= 2.75\times10^{-29}\ \mathrm{sr}
\end{align*}
\hfill(30 points)

The flux density is related to the total brightness by the relation:
\[ S_{\mathrm{1660\,MHz}} = B_{\mathrm{1660\,MHz}}\,\Omega \]
\hfill(20 points)

while at this frequency, the total brightness can be approximated from the
Rayleigh--Jeans formula:
\[ B_{\mathrm{1660\,MHz}} = 2kT_b\nu^2/c^2 \]
\hfill(20 points)

where $T_b$ is the brightness temperature. Then we have:
\begin{align*}
T_b &= S_{\mathrm{1660\,MHz}}c^2/2k\nu^2\Omega\\
T_b &= [3.5\times10^{-21}\times(3\times10^{10})^2]/
       [2\times(1.38\times10^{-16})\times(1660\times10^{6})^2
        \times(2.75\times10^{-29})]\\
    &= 1.5\times10^{26}\ \mathrm{K}
\end{align*}
\hfill(10 points)

\solhead{6.21}{21-cm HI galaxy survey}
% source id: 2017_TH_T4
\textbf{Case 1:} If we only consider the frequency range where HI spectral
line can be detected by the receiver of this radio telescope,
\[ z \approx \frac{f_0 - f}{f_0} \]
\hfill[1 Mark]

The lowest frequency, 1.32\,GHz, can be used to observe the highest
redshifted HI, at
\[ z = \frac{1.42 - 1.32}{1.42} = 0.0704 \]
\hfill[1 Mark]

\textbf{Alternative solution:} Above is the cosmological redshift formula in
frequency conventionally used in radio astronomy. However, if the
cosmological redshift in optical astronomy is used, student may use
\[ z \approx \frac{f_0 - f}{f} = \frac{1.42 - 1.32}{1.32} = 0.0758 \]
and the student should be awarded full mark for this part.

\textbf{Case 2:} Use the information given to calculate the redshift limit
set by the flux limit due to furthest distance that typical HI galaxy can be
detected by the telescope

for low $z$, non-relativistic redshift
\[ d = \frac{v_r}{H_0} = \frac{cz}{H_0} \]
\hfill[1 Mark]

The HI spectral-line flux density for a galaxy with HI luminosity $L$ and
line-width $\Delta f$ at distance $d$ is
\[ \frac{L}{4\pi d^2}\times\frac{1}{\Delta f}\
   (\mathrm{W\,m^{-2}\,Hz^{-1}}) \]
\hfill[2 Mark]

Setting this to the detection limit and solve correctly, and taking special
care of converting various units
\[ S = \frac{L H_0^2}{4\pi\Delta f c^2 z^2} \ge S_{\mathrm{lim}}
     = 0.5\times10^{-26} \]
\hfill[1 Mark]

Therefore,
\[ z \le \frac{H_0}{c}\sqrt{\frac{L}{4\pi\Delta f S_{\mathrm{lim}}}}, \]
\hfill[1 Mark]

Correctly substitute numerical values,
\[ z \le \frac{67.8\,(\mathrm{km^{-1}\,Mpc^{-1}})}
     {2.99\times10^{5}\,(\mathrm{km\ s^{-1}})\times
      3.08\times10^{22}\ (\mathrm{m\ Mpc^{-1}})}
   \sqrt{\frac{10^{28}\ \mathrm{W}}
     {4\pi\times10^{6}\,(\mathrm{Hz})\times0.5\times10^{-29}\
      (\mathrm{Wm^{-2}Hz^{-1}})}} \]
% sic: the numerator's units are printed "km^-1 Mpc^-1" in the source
\hfill[1 Mark]
\[ z \le 0.0929 \]
\hfill[1 Mark]

\textbf{Conclusion:} For typical HI galaxy, the highest redshift is limited
by the receiver's lowest frequency limit and we get $z_{\max} = 0.0704$ (or
0.0758 for alternative solution given above). \hfill[1 Marks]

\emph{The last part is marked only if student calculate both instrumental
factors. If students use more precise cosmological effects, they will be
marked accordingly.}

\solhead{6.22}{Supernova 1987A}
% source id: 2017_TH_T6
\begin{description}
\item[a)]
  \[ m - m_0 = -2.5\log_{10}\left(\frac{B}{B_0}\right) \]
  \hfill[1.0]
  \[ m - m_0 = 2.5\,\frac{t}{\tau}\log_{10} e \]
  \hfill[1.0]

  The days from 15\textsuperscript{th} May 1987 to 4\textsuperscript{th}
  February 1988 is 265 days \hfill[1.0]
  \[ 6 - 3 = 2.5\,\frac{265\ \mathrm{days}}{\tau}\log_{10} e \]
  \hfill[1.0]
  \[ \tau = 95.9\ \mathrm{days} \]
  \hfill[1.0]
\item[b)]
  Consider the energy transmission
  \[ B_e d^2 = T B_T D^2 \]
  \hfill[3.0]
  \begin{align*}
  \frac{B_e}{B_T} &= T\,\frac{D^2}{d^2}\\
  m - m_0 &= -2.5\log_{10}\left(\frac{B}{B_0}\right)
    \tag*{[1.0]}\\
  m_{\mathrm{lim}} &= m_e + 2.5\log_{10} T
    + 5\log_{10}\left(\frac{D}{d}\right)
    \tag*{[1.0]}\\
  m_{\mathrm{lim}} &= 6 + 2.5\log_{10} 0.7
    + 5\log_{10}\left(\frac{15.24\ \mathrm{cm}}{0.6\ \mathrm{cm}}\right)\\
  m_{\mathrm{lim}} &= 12.64
  \end{align*}
  and

  From $m - m_0 = 2.5\,\dfrac{t}{\tau}\log_{10} e$ \hfill[2.0]
  \[ 12.64 - 3 = 2.5\,\frac{t}{95.9\ \mathrm{days}}\log_{10} e \]
  \hfill[1.0]
  \[ t = 851.5\ \mathrm{days} \]
  \hfill[1.0]

  The last day that observers could have seen the supernova is on
  \textbf{12\textsuperscript{th} September 1989} \hfill[1.0]

  \emph{11\textsuperscript{th} -- 13\textsuperscript{th} September 1989 are
  acceptable}
\end{description}

\solhead{6.23}{Observe the Sun with FAST}
% source id: 2018_TH_T6
Rayleigh--Jeans law to calculate the thermal emission from the sun at 3\,GHz
states,
\[ B_\nu = \frac{2k_B T}{c^2}\,\nu^2 \]
which is the power emitted per unit emitting area, per steradian, per unit
frequency. Therefore, the solar luminosity at 3\,GHz should be:
\[ L_\nu = B_\nu \cdot 4\pi R_\odot^2 \]
\hfill(3 points)

At the distance of the earth (1\,AU), the monochromatic flux from the sun at
3\,GHz should be:
\[ f_\nu = \frac{L_\nu}{4\pi D^2} \]
Hence the energy flux that FAST will receive is:
\[ F_\nu = f_\nu \cdot \Delta\nu \cdot \pi^2\,\frac{d_c^2}{4} \]
\hfill(8 points)
% sic: the source prints a factor pi^2 here and in the next line; the
% collecting area of a dish of effective diameter d_c is pi d_c^2/4

And the total energy that the receiver will collect during 1 hour of
observation is:
\[ E_\odot = F_\nu\Delta t
   = \frac{2k_B T}{c^2}\,\nu^2\,\frac{R_\odot^2}{D^2}
     \cdot \pi^2\,\frac{d_c^2}{4}\,\Delta\nu\,\Delta t
   = 8.5\times10^{-5}\,J \]
\hfill(8 points)

Then we can calculate the work we need to turn over one piece of the answer
sheet (A4 paper). The mass of an A4 paper
($297\,\mathrm{mm}\times210\,\mathrm{mm}$) is:
\[ m = \rho \cdot L_1 L_2 \]
\hfill(2 point)

Therefore, the energy of turning it over should be around:
\[ E' = mg \cdot \frac{L_2}{2} \approx 5\times10^{-3}\,J \]
\hfill(3 points)

As a consequence, $E' > E_\odot$. \hfill(1 point)

\solhead{6.24}{Cosmic Microwave Background Oven}
% source id: 2019_TH_5
\begin{description}
\item[a)] If the body is spherical, its radius $r$ can be calculated as:
  \[ \text{(5.1)}\quad V = \frac{m}{\varrho} = \frac{4}{3}r^3\pi
     \;\rightarrow\; r = \left(\frac{3m}{4\pi\varrho}\right)^{1/3} \]
  The area of the body:
  \[ \text{(5.2)}\quad A = 4r^2\pi
     = 4\pi\left(\frac{3m}{4\pi\varrho}\right)^{2/3}
     = (4\pi)^{1/3}\left(\frac{3m}{\varrho}\right)^{2/3} \]
  \hfill(1 p)

  According to the Stefan-Boltzmann law, the total energy radiated per unit
  surface area across all wavelengths per unit time:
  \[ \text{(5.3)}\quad j = \sigma T_{\mathrm{CMB}}^4 \]
  \hfill(1 p)

  The CMB is isotropic, which means the absorbed radiation by unit surface
  area of the spherical body is the same independent of the surface's normal
  direction. \hfill(1 p)

  The total absorbed energy by the spherical body per unit time:
  \[ \text{(5.4)}\quad P = jA = \sigma T_{\mathrm{CMB}}^4
     \times (4\pi)^{1/3}\left(\frac{3m}{\varrho}\right)^{2/3} \]
  \hfill(1 p)

  Its numerical value:
  \[ P \approx \boxed{2.3\times10^{-6}\,\mathrm{W}} \]
  \hfill(1 p)
\item[b)] The mean energy of CMB photons:
  \[ \text{(5.5)}\quad \langle\varepsilon\rangle = 3kT_{\mathrm{CMB}}
     = 1.1\times10^{-22}\,\mathrm{J} \]
  \hfill(1 p)

  The number of CMB photons the body would absorb per second:
  \[ \text{(5.6)}\quad n = \frac{P}{\langle\varepsilon\rangle} \]
  \hfill(1 p)

  Its numerical value:
  \[ n = \boxed{2.1\times10^{16}\,\mathrm{s^{-1}}} \]
  \hfill(1 p)
\item[c)] Denote the time necessary for raising the temperature by
  $\Delta t$, so:
  % sic: the source says "by Delta t"; the temperature rise is Delta T
  \[ \text{(5.7)}\quad P\Delta t = Cm\Delta T
     \;\rightarrow\; \Delta t = \frac{Cm\Delta T}{P} \]
  \hfill(1 p)

  Its numerical value:
  \[ \Delta t \approx \boxed{1.1\times10^{11}\,\mathrm{s}}
     \approx 1\,253\,200\,\mathrm{d} \approx \boxed{3430\,\mathrm{yr}} \]
  \hfill(1 p)
\end{description}

\solhead{6.25}{Dyson Sphere}
% source id: 2022_TH_11
\begin{description}
\item[(a)] Heat absorbed must be fully emitted to maintain the thermal
  equilibrium
  \begin{align*}
  kL_\odot &= 4\pi R^2\epsilon\sigma T_{\mathrm{eq}}^4
    \tag*{2.0}\\
  \therefore\ T_{\mathrm{eq}} &=
    \sqrt[4]{\frac{kL_\odot}{4\pi R^2\epsilon\sigma}}
    \tag*{1.0}
  \end{align*}
\item[(b)] In this part, we try minimize the radius at the expanse of
  reaching the highest operational temperature for panels. Again
  \begin{align*}
  kL_\odot &= 4\pi R^2\epsilon\sigma T_{\max}^4\\
  R &= \sqrt{\frac{kL_\odot}{4\pi\epsilon\sigma T_{\max}^4}}
    \tag*{1.0}\\
  &\approx 1\times10^{11}\,\mathrm{m} = 0.665\,\mathrm{au}
    \tag*{2.0}
  \end{align*}
  Clearly, this distance is approximately 1.5 times smaller than Earth-Sun
  separation. Hence, Earth will stay \textbf{OUT}side the sphere. \hfill 1.0
\item[(c)] Power transmitted into usable energy
  \begin{align*}
  P &= \eta L_\odot\\
  P &= 7.65\times10^{25}\,\mathrm{W}
    \tag*{2.0}
  \end{align*}
\item[(d)] Energy harnessed in 1 second will be enough for
  \[ \tau = \frac{7.65\times10^{25}\,\mathrm{W}\times1\,\mathrm{s}}
       {17\times10^{12}\,\mathrm{W}}
     = 4.5\times10^{12}\,\mathrm{s} \approx 143\,000\,\mathrm{yr} \]
  \hfill 2.0
\item[(e)] New equilibrium temperature may be found as
  \begin{align*}
  kL_\odot\frac{\pi R_\oplus^2}{4\pi a_\oplus^2}
    &= 4\pi R_\oplus^2\sigma T_{\mathrm{new}}^4
    \tag*{2.0}\\
  T_{\mathrm{new}} &= \sqrt[4]{\frac{kL_\odot}{16\pi\sigma a_\oplus^2}}\\
  T_{\mathrm{new}} &\approx 206\,\mathrm{K}
    \tag*{2.0}\\
  \Delta T &= 288 - 206 = 82\,\mathrm{K}
    \tag*{1.0}
  \end{align*}
\item[(f)] By Kepler's Third Law
  \begin{align*}
  T^2 &= \frac{4\pi^2}{GM_\odot}R^3\\
  \therefore\ T &= (0.665)^{1.5}\,\mathrm{yr}\\
  T &\approx 0.542\,\mathrm{yr} = 198\,\mathrm{days}
    \tag*{3.0}
  \end{align*}
\item[(g)] We compare the two forces for given area $A$ of the solar panel
  \begin{align*}
  F_{\mathrm{RP}} &= \frac{I_{inc}A}{c} = \frac{L_\odot A}{4\pi R^2 c}
    \tag*{2.0}\\
  F_{\mathrm{Grav}} &= \frac{GM_\odot A\rho}{R^2}
    \tag*{1.0}\\
  \alpha &= \frac{F_{\mathrm{RP}}}{F_{\mathrm{Grav}}}
    = \frac{L_\odot}{4\pi cGM_\odot\rho}\\
  \alpha &\approx 7.65\times10^{-4}
    \tag*{2.0}
  \end{align*}
  At the first sight, this might seem like a negligible effect, but let us
  look at the change of the orbital period
  \[ (1-\alpha)F_{\mathrm{Grav}} = m\omega_1^2 R \]
  \hfill 3.0

  The radius should be the same, because our aim is to minimize the radius.
  \begin{align*}
  \omega_1 &= \sqrt{\frac{GM_\odot(1-\alpha)}{R^3}}\\
  T_1 &= 2\pi\sqrt{\frac{R^3}{GM_\odot}}\,(1-\alpha)^{-0.5}
    \tag*{1.0}\\
  \therefore\ \frac{T_1}{T} &= (1-\alpha)^{-0.5} \approx 1 + 0.5\alpha\\
  \Delta T &= 0.5\alpha T\\
  \Delta T &\approx 1.8\,\mathrm{h}
    \tag*{3.0}
  \end{align*}
  Therefore, \textbf{YES} this additional force has an easily detectable
  effect. \hfill 1.0
\item[(h)] First, we find the angle between entrance and exit points
  \begin{align*}
  r &= \frac{a}{1-\cos\theta}\\
  R - R\cos\theta &= a\\
  \cos\theta &= \frac{R-a}{R} = \frac{0.665-1}{0.665}
    \tag*{1.0}\\
  \cos\theta &= -0.5046
    \tag*{1.0}\\
  \theta_1 &= 2.099\,\mathrm{rad} = 120.2^\circ = 120^\circ14'\\
  \theta_2 &= 4.185\,\mathrm{rad} = 239.8^\circ = 239^\circ46'\\
  2\phi &= 2.086\,\mathrm{rad} = 119.5^\circ = 119^\circ32'
    \tag*{3.0}
  \end{align*}
  \ioaafig{2022_TH_11-s1.pdf}
  \ioaafig{2022_TH_11-s2.pdf}

  We know the that asteroid will stay inside the sphere for about
  $\tau_0 = 37$ days. During this time, sphere will rotate with an angle
  % sic: "We know the that asteroid" in source
  \begin{align*}
  \Delta\gamma &= \omega\tau_0 = \frac{2\pi\tau_0}{T}\\
  \Delta\gamma &\approx 1.17\,\mathrm{rad}
    \tag*{2.0}
  \end{align*}
  Hence, if the angular distance between holes is $\beta$, the asteroid to
  go through safely, the following condition must be satisfied
  \begin{align*}
  2\phi &= \Delta\gamma + \beta
    \tag*{5.0}\\
  \therefore\ \beta &= 2\phi - \Delta\gamma\\
  \beta &\approx 0.91\,\mathrm{rad} = 52.3^\circ
    \tag*{1.0}
  \end{align*}
\item[(i)] By Wien's law, the most probable wavelength radiated by the
  sphere would be
  \begin{align*}
  \lambda_1 &= \frac{b}{T_1}\\
  \lambda_2 &= \frac{b}{T_2}
    \tag*{2.0}\\
  \lambda_1' &\approx \lambda_1\left(1+\frac{dH_0}{c}\right)\\
  \lambda_2' &\approx \lambda_2\left(1+\frac{dH_0}{c}\right)
    \tag*{2.0}
  \end{align*}
  \[ \frac{b}{T_2}\left(1+\frac{dH_0}{c}\right) < \lambda_{obs}
     < \frac{b}{T_1}\left(1+\frac{dH_0}{c}\right) \]
  \hfill 1.0
\end{description}

\solhead{6.26}{Magnetic Field}
% source id: 2023_TH_Q2
\begin{description}
\item[(a)] In a magnetic field, a charged particle moves along a circular
  path defined by the equality of centrifugal and magnetic forces:
  \[ mv^2/r = evB, \]
  \hfill(1 point)

  where $m$ is the mass, $v$ velocity, $r$ radius of the circle, $e$ charge
  of the particle, and $B$ magnetic flux density.

  For circular motion, $v = 2\pi r/T$, therefore $T = 2\pi m/eB$. The
  charged particle moving in harmonic motion (i.e.\ along the circular path
  with constant velocity) emits a wave of wavelength
  $\lambda = cT = 2\pi mc/eB$ and thus $B = 2\pi mc/e\lambda$.
  \hfill(1 point)

  Substituting the numerical values, including the mass and charge of the
  electron:
  \[ B = (2\pi\times9.109\times10^{-31}\times2.998\times10^{8})/
     (1.602\times10^{-19}\times6\times10^{-7})
     \approx 2\times10^{4}\ \mathrm{T}. \]
  \hfill(1 point)

  Since $\lambda$ is given to 1 s.f., the correct answer is 20\,kT, however
  accept 17.9\,kT or 18\,kT. More than 3 s.f.\ in the final answer loses
  points.
\item[(b)] In the plasma, besides electrons, only protons will be present in
  large quantities; protons will emit energy at a wavelength larger in
  proportion to the mass ratio, i.e.\ $1836\times$ larger (anything within
  1800--2000$\times$ $\implies$ $\lambda = 1.08$--$1.20$\,mm is acceptable).
  \hfill(2 points)
\end{description}

\solhead{6.27}{Balmer Decrement}
% source id: 2025_TH_T04
\begin{description}
\item[(T04.1)] For H$\alpha$ substituting in equation for $\kappa(\lambda)$,
  we get
  \[ \kappa(\mathrm{H}\alpha)
     = 2.659\times\left(-1.857+\frac{1.040}{0.6563}\right) + 3.1 \]
  \[ \boxed{\kappa(\mathrm{H}\alpha) = 2.4} \]
  \hfill 1.5

  Similarly, for H$\beta$ we have
  \[ \kappa(\mathrm{H}\beta)
     = 2.659\times\left(-2.156+\frac{1.509}{0.4861}
       -\frac{0.198}{0.4861^2}+\frac{0.011}{0.4861^3}\right) + 3.1 \]
  \[ \boxed{\kappa(\mathrm{H}\beta) = 3.6} \]
  \hfill 1.5
\item[(T04.2)]
  \begin{align*}
  E(\mathrm{H}\beta-\mathrm{H}\alpha)
    &= A_{\mathrm{H}\beta} - A_{\mathrm{H}\alpha}
    \tag*{2.0}\\
    &= \left[\kappa(\mathrm{H}\beta)-\kappa(\mathrm{H}\alpha)\right]
       E(B-V) = 1.270\,E(B-V)
  \end{align*}
  \[ \boxed{\frac{E(\mathrm{H}\beta-\mathrm{H}\alpha)}{E(B-V)} = 1.2} \]
  \hfill 2.0
\item[(T04.3)] If $f_{\lambda,\mathrm{obs}}$ and $f_{\lambda,\mathrm{em}}$
  represent the observed and emitted fluxes at wavelength $\lambda$,
  respectively, then the extinction is given by the equation
  \[ A_\lambda = -2.5\log
     \left(\frac{f_{\lambda,\mathrm{obs}}}{f_{\lambda,\mathrm{em}}}\right) \]
  \hfill 1.0

  Given, $\dfrac{f_{\mathrm{H}\alpha,\mathrm{em}}}
    {f_{\mathrm{H}\beta,\mathrm{em}}} = 2.86$
  \begin{align*}
  E(\mathrm{H}\beta-\mathrm{H}\alpha)
    &= A_{\mathrm{H}\beta} - A_{\mathrm{H}\alpha}\\
    &= -2.5\log\left(
       \frac{f_{\mathrm{H}\beta,\mathrm{obs}}}
            {f_{\mathrm{H}\alpha,\mathrm{obs}}}\times
       \frac{f_{\mathrm{H}\alpha,\mathrm{em}}}
            {f_{\mathrm{H}\beta,\mathrm{em}}}\right)\\
    &= -2.5\log\left(\frac{1.06\times10^{-15}}{6.80\times10^{-15}}
       \times2.86\right)\\
    &= 0.88
    \tag*{1.5}
  \end{align*}
  From extinction curve, we have
  $E(\mathrm{H}\beta-\mathrm{H}\alpha) = 1.2\,E(B-V)$.\\
  Equating these two, we obtain
  \[ E(B-V) = \frac{0.88}{1.2} = 0.73 \]
  \hfill 0.5

  This gives,
  \[ A_{\mathrm{H}\alpha} = \kappa(\mathrm{H}\alpha)\,E(B-V)
     = 2.4\times0.73 \]
  \[ \boxed{A_{\mathrm{H}\alpha} = 1.8\,\mathrm{mag}} \]
  \hfill 1.5
  \[ A_{\mathrm{H}\beta} = \kappa(\mathrm{H}\beta)\,E(B-V)
     = 3.6\times0.73 \]
  \[ \boxed{A_{\mathrm{H}\beta} = 2.6\,\mathrm{mag}} \]
  \hfill 1.5
\item[(T04.4)]
  \[ A_V = 3.1\times E(B-V) = 3.1\times0.69 \]
  % sic: the source substitutes 0.69 here although E(B-V) = 0.73 was
  % obtained above (full credit was given for E(B-V) between 0.69 and 0.73)
  \[ \boxed{A_V = 2.3\,\mathrm{mag}} \]
  \hfill 1.0
  \[ m_{\mathrm{V0}} = m_V - A_V = 11.315 - 2.3 \]
  \[ \boxed{m_{\mathrm{V0}} = 9.0\,\mathrm{mag}} \]
  \hfill 1.0
\end{description}

\solhead{6.28}{}
% source id: 2013_DA_D3
\textbf{Answer:} The human eye can easily recognize differences of the order
of $\Delta m_V = 0.1$ -- $0.2$\,mag. The student should first align the
photograph with the chart, which has a different scale and orientation. Then
the magnitudes of the stars can easily be recognized. Note: not all arrowed
stars can be found in the photograph. (0.5 Point for each star that has been
correctly identified and 0.5 point for every correct magnitude within
$\pm0.3$ magnitude \hfill(18 Points)

\begin{center}
\begin{tabular}{|c|c|c|c|c|c|c|c|c|c|}
\hline
\textbf{Star} & 3 & 4 & 5 & 6 & 7 & 8 & 9 & 10 & 11\\
\hline
\textbf{m} & 7.3 -- 7.9 & 5.3 -- 5.9 & 5.2 -- 5.8 & 4.6 -- 5.2
  & 5.2 -- 5.8 & 4.5 -- 5.1 & 4.3 -- 4.9 & 4.9 5.5 % sic: dash missing
  & 5.3 -- 5.9\\
\hline
\end{tabular}
\end{center}

Stars 1, 2 and 12 are outside the boundaries of the photograph
\ioaafig{2013_DA_D3-s1.pdf}

\solhead{6.29}{Observer on an extrasolar planet}
% source id: 2014_DA_D3
% Official solution ("Problem 3 B. Marking scheme -- Observer on an extrasolar
% Note: the official solution names the asterism "Big Dipper" (and keeps some
% Romanian words: Planeta, Soare, Pamant); the official English statement --
% and the book's problem text -- uses the Mizar system instead. Transcribed
% as printed.

\begin{description}
\item[a)] The magnitude of the \textbf{SUN as seen from the observer on the
  Sirrius planet} % sic: "Sirrius" in source
  \[ \log\frac{E_{\mathrm{SUN},EARTH}}{E_{\mathrm{SUN},PLANET}}
     = -0.4\left(m_{\mathrm{SUN},EARTH} - m_{\mathrm{SUN},PLANET}\right), \]
  Where: $E$ represents the apparent brightness of the body at the
  corresponding distance to the observer.
  \begin{align*}
  \log\frac{\dfrac{L_{\mathrm{S}}}{4\pi d_{\mathrm{S},E}^2}}
           {\dfrac{L_{\mathrm{S}}}{4\pi d_{\mathrm{S,P}}^2}}
    &= -0.4\left(m_{\mathrm{S},E} - m_{\mathrm{S,P}}\right);\\
  d_{\mathrm{S},E} = 1\,\mathrm{A}U;\
  d_{\mathrm{S,P}} &= 2.6\,\mathrm{pc}
    = 2.6\cdot206265\ \mathrm{A}U;\\
  m_{\mathrm{S},E} &= -26.78^{\mathrm{m}};\\
  m_{\mathrm{S,P}} &= 1.87^{\mathrm{m}},
  \end{align*}
  Apparent magnitude of Sun seen from the observer on one of the Sirius
  planet.
\item[b)] \textbf{Magnitude of Sirius, seen from the observer from its
  planet}:
  \begin{align*}
  \log\frac{E_{\mathrm{Sirius,Planeta}}}{E_{\mathrm{Sun,Planeta}}}
    &= -0.4\left(m_{\mathrm{Sirius,Planeta}}
       - m_{\mathrm{Sun,Planeta}}\right),\\
  d_{\mathrm{Sun,Planeta}} &= 2.6\,\mathrm{pc}
    = 2.6\cdot206265\ A\mathrm{U};\\
  d_{\mathrm{Sirius,Planeta}} &= 10\ A\mathrm{U};\\
  L_{\mathrm{Sirius}} &= 25\cdot L_{SUN};\\
  m_{\mathrm{Sun,Planeta}} &= 1.87^{\mathrm{m}};\\
  m_{\mathrm{Sirius,Planeta}} &= -25.255^{\mathrm{m}},
  \end{align*}
\item[c)] The luminosity of \textbf{Big Dipper}:
  \[ L_{\mathrm{Big\ Dipper}} = L_1 + L_2 + L_3 + L_4 + L_5 + L_6 + L_7, \]
  where $L_1$, $L_2$, \dots, $L_7$ are the luminosties % sic: "luminosties"
  of the seven stars of the asterism:
  \begin{align*}
  \log\frac{E_{\mathrm{S}un,Earth}}{E_{1,Earth}}
    &= -0.4\left(m_{\mathrm{SUN},Earth} - m_{1,Earth}\right),\\
  \log\frac{\dfrac{L_{\mathrm{S}un}}{4\pi d_{\mathrm{S}un,Earth}^2}}
           {\dfrac{L_1}{4\pi d_{1,Earth}^2}}
    &= -0.4\left(m_{\mathrm{S}un,Earth} - m_{1,Earth}\right);\\
  d_{1,Earth} &= 124\cdot63240\ A\mathrm{U};\\
  m_{\mathrm{S}un,Earth} = -26.78^{\mathrm{m}};\
  m_{1,Earth} &= 1.8^{\mathrm{m}};\\
  L_1 &\approx 225\cdot L_{\mathrm{S}},
  \end{align*}
  The luminosity of star 1 (Dubhe) from the asterism Big Dipper.

  Symilarly for the rest of stars in the asterism: % sic: "Symilarlyfor"
  \begin{align*}
  L_2 &\approx 50\cdot L_{\mathrm{S}};\\
  L_3 &\approx 60\cdot L_{\mathrm{S}};\\
  L_4 &\approx 0.134\cdot L_{\mathrm{S}};
    % sic: 0,134 as printed; the data of the problem (m = 3,3 at 58 ly)
    % give L_4 of order 13 L_S
  \\
  L_5 &\approx 52\cdot L_{\mathrm{S}};\\
  L_6 &\approx 4\cdot L_{\mathrm{S}};\\
  L_7 &\approx 1.37\cdot L_{\mathrm{S}}.
  \end{align*}
  The luminosity of the asterism Big Dipper is:
  \begin{align*}
  L_{\mathrm{Big\ Dipper}} &= L_1 + L_2 + L_3 + L_4 + L_5 + L_6 + L_7;\\
  L_{\mathrm{Big\ Dipper}} &= (225 + 50 + 60 + 0.134 + 52 + 4
    + 1.37)L_{Sun};\\
  L_{\mathrm{Big\ Dipper}} &= 392.5\cdot L_{\mathrm{S}un}.
  \end{align*}
\item[d)] \textbf{The distance between Big Dipper and Earth}
  $d_{\mathrm{Big\ Dipper},Earth}$:
  \begin{align*}
  \log\frac{E_{\mathrm{Big\ Dipper},Earth}}{E_{\mathrm{Soare},Earth}}
    &= -0.4\left(m_{\mathrm{Big\ Dipper},Earth}
       - m_{\mathrm{Soare},Earth}\right),\\
  L_{\mathrm{Big\ Dipper}} &= 392.5\cdot L_{\mathrm{S}};\\
  d_{\mathrm{Big\ Dipper},Earth} &= 573\cdot10^{3}\cdot\sqrt{392.5}
    \cdot d_{\mathrm{Soare},Earth};\\
  d_{\mathrm{Big\ Dipper},Earth} &\approx 11.3\cdot10^{6}\ AU.
  \end{align*}
\item[e)] Geocentric angular distance $\Delta\theta$, \textbf{between
  Big Dipper and Sirius}. % sic: source prints "between dintre Big Dipper"

  The equatorial coordinates of asterism Big Dipper $(\sigma_1)$ and Sirius
  $(\sigma_2)$ located on the sphere of geocentric celestial sphere are:
  \begin{align*}
  \alpha_{\mathrm{Big\ Dipper}} &= \alpha_1
    = 13^{\mathrm{h}}23^{\mathrm{min}}55.5^{\mathrm{s}};\
  \delta_{\mathrm{Big\ Dipper}} = \delta_1 = 54^\circ55'31'';\\
  \alpha_{\mathrm{Sirius}} &= \alpha_2 = 6^{\mathrm{h}}45^{\mathrm{min}};\
  \delta_{\mathrm{Sirius}} = \delta_2 = -16^\circ43'.
  \end{align*}
  \ioaafig{2014_DA_D3-s1.pdf}

  In the spherical dreptunghic triangle ABC: % sic: "dreptunghic"
  % (Romanian: right-angled)
  \[ \sphericalangle\left(\sigma_1\,\mathrm{O}\,\sigma_2\right)
     = \Delta\theta. \]
  Also
  \begin{align*}
  \sphericalangle\left(\gamma\,\mathrm{O}\,\sigma_{01}\right)
    &= \alpha_{\mathrm{Big\ Dipper}} = \alpha_1;\
  \sphericalangle\left(\sigma_{01}\,\mathrm{OB}\right)
    = \delta_{\mathrm{Big\ Dipper}} = \delta_1;\\
  \sphericalangle\left(\gamma\,\mathrm{O}\,\sigma_{02}\right)
    &= \alpha_{\mathrm{Sirius}} = \alpha_2;\
  \sphericalangle\left(\sigma_{02}\,\mathrm{OC}\right)
    = \delta_{\mathrm{Sirius}} = \delta_2;\\
  \sphericalangle\left(\sigma_{01}\,\mathrm{O}\,\sigma_{02}\right)
    &= \sphericalangle\left(\mathrm{AOC}\right)
    = \alpha_1 - \alpha_2 = \Delta\alpha;\\
  \sphericalangle\left(\sigma_1\,\mathrm{OA}\right)
    &= \sphericalangle\left(\sigma_{01}\,\mathrm{O}\,\sigma_1\right)
     + \sphericalangle\left(\sigma_{01}\,\mathrm{OA}\right)
    = \delta_{\mathrm{Big\ Dipper}} - \delta_{\mathrm{Sirius}}
    = \Delta\delta;\\
  \sphericalangle\left(\sigma_1\,\mathrm{O}\,\sigma_2\right)
    &= \Delta\theta.
  \end{align*}
  Results
  \begin{align*}
  \cos a &= \cos b\cdot\cos c + \sin b\cdot\sin c\cdot\cos A;\\
  A &= 90^{0};\\
  \cos a &= \cos b\cdot\cos c;\\
  a = \Delta\theta;\ b &= \Delta\alpha;\quad c = \Delta\delta;\\
  \cos\Delta\theta &= \cos\Delta\alpha\cdot\cos\Delta\delta;\\
  \Delta\alpha = \alpha_1 - \alpha_2
    &= 13^{\mathrm{h}}23^{\mathrm{min}}55.5^{\mathrm{s}}
     - 6^{\mathrm{h}}45^{\mathrm{min}}
     = 6^{\mathrm{h}}38^{\mathrm{min}}55.5' \approx 100^\circ;
    % sic: 55,5' for seconds as printed
  \\
  \Delta\delta = \delta_1 - \delta_2
    &= 54^\circ55'31'' + 16^\circ43' = 71^\circ38'31''
    \approx 72.4^\circ;\\
  \cos\Delta\theta &= \cos\left(100^\circ\right)
    \cdot\cos\left(72.4^\circ\right)
    = \left(-0.173648177\right)\cdot\left(0.30236989\right)
    = -0.05250598;\\
  \Delta\theta &\approx 90.3^\circ \approx 90^\circ,
  \end{align*}
\item[f)] The distance between \textbf{Big Dipper and the observer on sirius
  planet of Sirius}, % sic: as printed
  $d_{\mathrm{Big\ Dipper,Planeta}}$:
  \begin{align*}
  d_{\mathrm{Big\ Dipper,Planeta}}
    &\approx d_{\mathrm{Big\ Dipper,Sirius}};\\
  d_{\mathrm{Sirius},Earth} &\approx d_{\mathrm{Sirius,Soare}}
    \approx d_{\mathrm{Soare,Planeta}};\\
  d_{\mathrm{Sirius,Pamant}} &= 2.6\,\mathrm{pc};\\
  d_{\mathrm{Soare,Planeta}} &= 2.6\,\mathrm{pc}
    = 2.6\cdot206265\ \mathrm{UA} = d_{Earth};\\
  d_{\mathrm{Big\ Dipper},Earth} &\approx 11.3\cdot10^{6}\ AU;
  \end{align*}
  \ioaafig{2014_DA_D3-s2.pdf}
  \[ d_{\mathrm{Big\ Dipper,Sirius}}
     = \sqrt{d_{\mathrm{Big\ Dipper,Pamant}}^2
       + d_{\mathrm{Sirius},Eaarth}^2} % sic: "Eaarth" in source
     = d_{\mathrm{Big\ Dipper,Planeta}}
     \approx 11.312\cdot10^{6}\ \mathrm{UA}. \]
\item[g)] The magnitude of \textbf{Big Dipper for an observer from the
  observer on planet of Sirius}: % sic: as printed
  \begin{align*}
  \log\frac{E_{\mathrm{Big\ Dipper,Planeta}}}
           {E_{\mathrm{Big\ Dipper},Earth}}
    &= -0.4\left(m_{\mathrm{Big\ Dipper,Planeta}}
       - m_{\mathrm{Big\ Dipper},Earth}\right),\\
  m_{\mathrm{Big\ Dipper},Earth} \approx 2^{\mathrm{m}};\
  d_{\mathrm{Big\ Dipper},Earth} &\approx 11.3\cdot10^{6}\ A\mathrm{U};\\
  d_{\mathrm{Big\ Dipper,Sirius}} = d_{\mathrm{Big\ Dipper,Planeta}}
    &\approx 11.312\cdot10^{6}\ \mathrm{A}U;\\
  m_{\mathrm{Big\ Dipper,Planeta}}
    &= 2^{\mathrm{m}} + 5^{\mathrm{m}}\cdot
       \log\frac{11.312\cdot10^{6}\ \mathrm{AU}}
                {11.3\cdot10^{6}\ \mathrm{AU}};\\
  m_{\mathrm{Big\ Dipper,Planeta}}
    &= 2^{\mathrm{m}} + 5^{\mathrm{m}}\cdot
       \log\frac{11.312}{11.3};\\
  m_{\mathrm{Big\ Dipper,Planeta}} &\approx 2.0023^{\mathrm{m}},
  \end{align*}
\end{description}

\solhead{6.30}{Interstellar extinction / IC 342}
% source id: 2015_DA_D2
\begin{description}
\item[a)] Using the data from tables 5 to 9 and recalling that MK class II
  corresponds to giant stars while MK classes Ia and Iab correspond to
  supergiants, we easily obtain tables 10 to 12 and hence figures containing
  the plots of $E_{\mathrm{X-V}}/E_{\mathrm{B-V}}$ against $1/\lambda_X$ for
  both stars.

  \begin{center}
  Table 10 (10 points)

  \begin{tabular}{|l|c|c|c|c|c|c|c|c|c|c|}
  \hline
  Star & $\frac{B-V}{\mathrm{mag}}$ & $\frac{V-V}{\mathrm{mag}}$
    & $\frac{R-V}{\mathrm{mag}}$ & $\frac{I-V}{\mathrm{mag}}$
    & $\frac{J-V}{\mathrm{mag}}$ & $\frac{H-V}{\mathrm{mag}}$
    & $\frac{K-V}{\mathrm{mag}}$ & $\frac{L-V}{\mathrm{mag}}$
    & $\frac{M-V}{\mathrm{mag}}$ & $\frac{N-V}{\mathrm{mag}}$\\
  \hline
  HD 4817 & 1.9 & 0 & $-1.45$ & $-2.54$ & $-3.42$ & $-4.32$ & $-4.64$
    & $-4.86$ & $-4.59$ & --\\
  \hline
  HD 11092 & 2.09 & 0 & -- & -- & $-3.47$ & $-4.43$ & $-4.94$ & $-5.16$
    & $-4.92$ & $-5.13$\\
  \hline
  \end{tabular}
  \end{center}

  \begin{center}
  Table 11 (10 points)

  \begin{tabular}{|l|c|c|c|c|c|c|c|c|c|c|}
  \hline
  Star & $\frac{(B-V)_0}{\mathrm{mag}}$ & $\frac{(V-V)_0}{\mathrm{mag}}$
    & $\frac{(R-V)_0}{\mathrm{mag}}$ & $\frac{(I-V)_0}{\mathrm{mag}}$
    & $\frac{(J-V)_0}{\mathrm{mag}}$ & $\frac{(H-V)_0}{\mathrm{mag}}$
    & $\frac{(K-V)_0}{\mathrm{mag}}$ & $\frac{(L-V)_0}{\mathrm{mag}}$
    & $\frac{(M-V)_0}{\mathrm{mag}}$ & $\frac{(N-V)_0}{\mathrm{mag}}$\\
  \hline
  HD 4817 & 1.42 & 0 & $-1.13$ & $-1.96$ & $-2.41$ & $-3.14$ & $-3.25$
    & $-3.39$ & $-3.25$ & $-3.63$\\
  \hline
  HD 11092 & 1.42 & 0 & $-0.96$ & $-1.61$ & $-2.16$ & $-2.77$ & $-3.05$
    & $-3.22$ & $-3.08$ & $-3.02$\\
  \hline
  \end{tabular}
  \end{center}

  \begin{center}
  Table 12 (10 points)

  \begin{tabular}{|l|c|c|c|c|c|c|c|c|c|c|}
  \hline
  Star & $\frac{E_{\mathrm{B-V}}}{E_{\mathrm{B-V}}}$
    & $\frac{E_{\mathrm{V-V}}}{E_{\mathrm{B-V}}}$
    & $\frac{E_{\mathrm{R-V}}}{E_{\mathrm{B-V}}}$
    & $\frac{E_{\mathrm{I-V}}}{E_{\mathrm{B-V}}}$
    & $\frac{E_{\mathrm{J-V}}}{E_{\mathrm{B-V}}}$
    & $\frac{E_{\mathrm{H-V}}}{E_{\mathrm{B-V}}}$
    & $\frac{E_{\mathrm{K-V}}}{E_{\mathrm{B-V}}}$
    & $\frac{E_{\mathrm{L-V}}}{E_{\mathrm{B-V}}}$
    & $\frac{E_{\mathrm{M-V}}}{E_{\mathrm{B-V}}}$
    & $\frac{E_{\mathrm{N-V}}}{E_{\mathrm{B-V}}}$\\
  \hline
  HD 4817 & 1.00 & 0.00 & $-0.67$ & $-1.21$ & $-2.10$ & $-2.46$ & $-2.90$
    & $-3.06$ & $-2.79$ & --\\
  \hline
  HD 11092 & 1.00 & 0.00 & -- & -- & $-1.96$ & $-2.48$ & $-2.82$ & $-2.90$
    & $-2.75$ & $-3.15$\\
  \hline
  \end{tabular}
  \end{center}

  \ioaafig{2015_DA_D2-s1.pdf}
  \ioaafig{2015_DA_D2-s2.pdf}
\item[b)] We have that
  $E_{\lambda-\mathrm{V}} = A_\lambda - A_{\mathrm{V}}
  \rightarrow -A_{\mathrm{V}}$ as $\lambda\rightarrow\infty$, so the plot of
  $E_{\mathrm{X-V}}/E_{\mathrm{B-V}}$ against $1/\lambda_X$ intersects the
  vertical axis at
  \[ \frac{-A_{\mathrm{V}}}{E_{\mathrm{B-V}}} = -R_{\mathrm{V}}. \]
  \hfill(10 points)

  Hence $R_{\mathrm{V}}$ can be read off as minus the intersection of the
  plot with vertical axis. For the two stars we get Table 13 (the
  intersection points are obtained by fitting a curve to guide the eye,
  noting that as $1/\lambda\rightarrow0$, the curve becomes flat, as hinted
  at above).

  \begin{center}
  Table 13 (10 points)

  \begin{tabular}{|l|c|}
  \hline
  Star & $R_{\mathrm{V}}$\\
  \hline
  HD 4817 & 3.1\\
  \hline
  HD 11092 & 3.0\\
  \hline
  \end{tabular}
  \end{center}

  Next,
  \[ R_{\mathrm{r,r-i}} = \frac{A_{\mathrm{r}}}{E_{\mathrm{r-i}}}
     = \left(\frac{E_{\mathrm{r-V}}}{E_{\mathrm{B-V}}}
       + \frac{A_{\mathrm{V}}}{E_{\mathrm{B-V}}}\right)
       \frac{E_{\mathrm{B-V}}}{E_{\mathrm{r-i}}}
     = \frac{\dfrac{E_{\mathrm{r-V}}}{E_{\mathrm{B-V}}} + R_{\mathrm{V}}}
       {\dfrac{E_{\mathrm{r-V}}}{E_{\mathrm{B-V}}}
        - \dfrac{E_{\mathrm{i-V}}}{E_{\mathrm{B-V}}}}, \]
  where $E_{\mathrm{r-V}}/E_{\mathrm{B-V}}$ and
  $E_{\mathrm{i-V}}/E_{\mathrm{B-V}}$ can be read off from the plot (again,
  by fitting a curve to guide the eye). Hence we get Table 14.

  \begin{center}
  Table 14 (10 points)

  \begin{tabular}{|l|c|}
  \hline
  Star & $R_{\mathrm{r}}$\\
  \hline
  HD 4817 & 3.7\\
  \hline
  HD 11092 & 3.6\\
  \hline
  \end{tabular}
  \end{center}

  Hence we take the expected values to be $R_{\mathrm{V}} \approx 3.1$ and
  $R_{\mathrm{r}} \approx 3.7$. Note that the first value agrees with the
  widely used ratio of the total to selective extinction in filters B and V.
\item[c)] First let us find the apparent distance moduli $\mu_{\mathrm{r,i}}$
  in filters r and i. Reading off the fitted values e.g.\ at
  $\log(P/\mathrm{day}) = 1.6$ from figures 2 and 3 and substituting into
  the period-luminosity relations, we find $\mu_{\mathrm{r}} = 29.0$\,mag
  and $\mu_{\mathrm{i}} = 28.6$\,mag, so
  $E_{\mathrm{r-i}} = A_{\mathrm{r}} - A_{\mathrm{i}}
  = \mu_{\mathrm{r}} - \mu_{\mathrm{i}} = 0.4$\,mag and so
  $A_{\mathrm{r}} \approx 3.7E_{\mathrm{r-i}} = 1.5$\,mag. Hence the
  unreddened distance modulus is
  $\mu_0 = \mu_{\mathrm{r}} - A_{\mathrm{r}} = 27.5$\,mag and so we estimate
  the distance to IC 342 to be 3.2\,Mpc. \hfill(20 points)
\end{description}

\solhead{6.31}{Naked-eye magnitude estimates near Cepheus}
% source id: 2009_OBS-SKY_2
% Official solution: "Solution 2" of the 2009 observation round
% ($IOAA_PAPERS_DIR/observation/2009_OBS_sol.pdf, pp. 87-88).

\textbf{Part 1:}

\begin{center}
Table 2

\begin{tabular}{|c|c|}
\hline
\multicolumn{2}{|c|}{\textbf{Angular Distance}}\\
\hline
\textbf{Stars Name} & \textbf{Angular Distance (Degree)}\\
\hline
1 (Errai) and 2 (Alfirk) & $11^\circ\!:09':10'' \sim 11^\circ$\\
\hline
1 (Errai) and 3 (Alderamin) & $18^\circ\!:36':50'' \sim 19^\circ$\\
\hline
\end{tabular}
\end{center}

\textbf{Part 2:}

\begin{center}
Table 4

\begin{tabular}{|c|c|c|}
\hline
\multicolumn{3}{|c|}{\textbf{Magnitude Estimation}}\\
\hline
\textbf{Number} & \textbf{Star Name} & \textbf{Visible Magnitude}\\
\hline
2 & Alfirk & 3.2\\
\hline
3 & Alderamin & 2.4\\
\hline
\end{tabular}
\end{center}

\solhead{6.32}{Magnitude of a missing star in NGC 6231}
% source id: 2012_OBS-TEL_2
% Official solution: item 2 of "Observational exam -- 1st night -- Solutions"
% ($IOAA_PAPERS_DIR/observation/2012_OBS_sol.pdf, p.1).

Exact answer is magnitude 7.9 but we will assume values between 7.7 -- 8.1 as
correct (0.2 magnitude tolerance). Values between 7.5 -- 7.7 or 8.1 -- 8.3
will be considered as 70\% correct. (5 minutes)
